%% LyX 2.4.0-alpha3 created this file.  For more info, see https://www.lyx.org/.
%% Do not edit unless you really know what you are doing.
\documentclass[english]{article}
\usepackage{lmodern}
\renewcommand{\sfdefault}{lmss}
\renewcommand{\ttdefault}{lmtt}
\usepackage[T1]{fontenc}
\usepackage{textcomp}
\usepackage[latin9]{inputenc}
\synctex=-1
\usepackage{color}
\definecolor{shadecolor}{rgb}{1, 0.667969, 0.5}
\usepackage{babel}
\usepackage{cprotect}
\usepackage{wrapfig}
\usepackage{booktabs}
\usepackage{calc}
\usepackage{framed}
\usepackage{url}
\usepackage{amsmath}
\usepackage{amssymb}
\usepackage{esint}
\usepackage[unicode=true,pdfusetitle,
 bookmarks=true,bookmarksnumbered=true,bookmarksopen=true,bookmarksopenlevel=2,
 breaklinks=true,pdfborder={0 0 1},backref=section,colorlinks=true,pdfpagemode=FullScreen]
 {hyperref}
\hypersetup{
 pdfborderstyle=}

\makeatletter

%%%%%%%%%%%%%%%%%%%%%%%%%%%%%% LyX specific LaTeX commands.
\newcommand{\noun}[1]{\textsc{#1}}
%% Because html converters don't know tabularnewline
\providecommand{\tabularnewline}{\\}

%%%%%%%%%%%%%%%%%%%%%%%%%%%%%% Textclass specific LaTeX commands.
\newenvironment{centred}%
  {\begin{center}\baselineskip=13pt\parskip=1pt }{\end{center}}
\newenvironment{example}%
{\begin{center}\ttfamily}{\end{center}}
\newenvironment{lyxcode}
	{\par\begin{list}{}{
		\setlength{\rightmargin}{\leftmargin}
		\setlength{\listparindent}{0pt}% needed for AMS classes
		\raggedright
		\setlength{\itemsep}{0pt}
		\setlength{\parsep}{0pt}
		\normalfont\ttfamily}%
	 \item[]}
	{\end{list}}

%%%%%%%%%%%%%%%%%%%%%%%%%%%%%% User specified LaTeX commands.
\usepackage[def-file=diffcoeff-doc]{diffcoeff}[=v4]

\makeatother

\begin{document}
\title{\texttt{diffcoeff}~\\
 a \LaTeX{} package to ease\texttt{ }the~\\
 writing of differential coefficients \\
 Version 4.2}
\author{Andrew Parsloe\\
{\small (ajparsloe@gmail.com)}}
\maketitle
\begin{abstract}
\noindent This package eases the consistent writing of ordinary, partial
and other derivatives of arbitrary (algebraic or numeric) order. For
mixed partial derivatives, the total order of differentiation is calculated
by the package. Optional arguments allow for points of evaluation
(ordinary derivatives), or variables held constant (partial derivatives),
and the placement of the differentiand in numerator or appended. The
package uses \texttt{xtemplate}, allowing systematic fine-tuning of
the display and generation and use of variant forms, including derivatives
built from $D$, $\Delta$ or $\delta$. A command for differentials
ensures the $\dl x$ (or $\dl.up.x$) used in integrals is consistent
with the form used in derivatives. 
\end{abstract}
\noindent %
\noindent\begin{minipage}[t]{1\columnwidth}%
\begin{shaded}%

\subsubsection*{Version 4.2/version 5.0 incompatibility}

Version 4.2 of \verb`diffcoeff` is issued at the same time as version
5.0. The latter is a complete rewrite of the package with many changes
to the user interface, arising above all from separating the fraction
types into separate commands (\verb`\diff`, \verb`\diffp` for upright
fractions; \verb`\difs`, \verb`\difsp` for slash fractions; and
\verb`\difc`, \verb`\difcp` for a new, compact style of derivative).
Although some elements of the earlier version are carried over, most
documents that use version 4 will fail to compile properly with version
5. To enable older documents to continue to compile, the \verb`\usepackage`
statement in the preamble needs a trailing optional argument added
to it:
\begin{verbatim}
    \usepackage[<options>]{diffcoeff}[=v4]
\end{verbatim}
The trailing argument \verb`[=v4]` (no spaces!) ensures \verb`diffcoeff4.sty`
rather than \verb`diffcoeff.sty` is used to compile the document.
Otherwise, version 4.2 makes no visible change to \verb`diffcoeff`
version 4.1 other than to documentation. (It also tidies a line of
code for the differential command \verb`\dl`.)\end{shaded}%
\end{minipage}

\tableofcontents{}

\section{Introduction}

The \LaTeX{} package \texttt{diffcoeff} is written in the expl3 language
of \LaTeX 3\texttt{ }and requires the bundles \texttt{l3kernel} which
is now part of standard \LaTeX{} (since early 2020) and \texttt{l3packages}
(the latter for the \texttt{xparse}, \texttt{l3keys2e} and especially
\texttt{xtemplate} packages\texttt{)}. The package is loaded by entering 
\begin{lyxcode}
\textbackslash usepackage{[}<options>{]}\{diffcoeff\}{[}=v4{]}
\end{lyxcode}
in the preamble of your document, where \verb`<options>` are package
options input by the user; see below. The trailing optional argument,
\verb`[=v4]`, is necessary to ensure your document compiles with
version 4 of \verb`diffcoeff` rather than the new version 5 in which
there are many changes to the user interface.

\subsubsection*{Note on terminology}

I refer throughout to the quantity or function being differentiated
as the \emph{differentiand} or \emph{derivand }(in line with \emph{integrand},
\emph{operand}, etc.).

\subsubsection*{New in version 4.0}

(See �\ref{sec:Version-comparison} for a comparison of versions.)
Version 4.0
\begin{enumerate}
\item enables the insertion of a small space before the differentiand, either
as the default behaviour (package option \verb`spaced`) or at explicit
request (argument \verb`!` of the \verb`\diff` command); see �\ref{subsec:Spacing-before-derivand};
\item prevents the ligature $df$ that marred previous versions; this is
now rendered $\dl f$;
\item offers the document command \verb`\dlp` for the \emph{partial} differential;
see �\ref{subsec:Partial-differential};
\item no longer accepts \emph{the braced form} of the\emph{ }trailing optional
argument specifying a point of evaluation or (for partial derivatives)
variables held constant. This was a relic from version 1 of \texttt{diffcoeff},
and has been deprecated since version 2.
\end{enumerate}

\subsubsection*{New in version 4.1}
\begin{enumerate}
\item changes the code handling the \verb`long-var-wrap` setting to avoid
a problem with a change in \verb`xtemplate` (see \url{https://tex.stackexchange.com/questions/646976/error-with-package-diffcoeff})\footnote{Based on a suggestion from \noun{Joseph Wright}.}; 
\item prevents (hitherto overlooked) ligatures in the \emph{denominators}
of differential coefficients.
\end{enumerate}

\subsubsection*{New in version 4.2}
\begin{enumerate}
\item tidies a line of code in the definition of the differential (has no
visible effect);
\item amends documentation in light of the issue of version 5 and the resulting
interface changes.
\end{enumerate}

\subsection{Package options}

\label{subsec:Package-options}There are three package options, which
are entered in a comma-separated list in the optional argument of
the \texttt{\textbackslash usepackage} command. 
\begin{enumerate}
\item The first is a switch, \texttt{ISO}, which turns on formatting conforming
to ISO recommendations, principally upright `d's: 
\begin{lyxcode}
\textbackslash usepackage{[}ISO{]}\{diffcoeff\}{[}=v4{]}
\end{lyxcode}
The full effect of this option is disscussed in �\ref{sec:Changing-defaults}. 
\item The second is another switch, \texttt{spaced}, and \emph{is new with
version 4 }of \texttt{diffcoeff}. This switch ensures a small space
is  inserted before the differentiand:
\begin{lyxcode}
\textbackslash usepackage{[}spaced{]}\{diffcoeff\}{[}=v4{]}
\end{lyxcode}
This is discussed in �\ref{subsec:Spacing-before-derivand}.
\item The third requires a filename \texttt{<filename>} for a file \texttt{<filename>.def}
containing definitions of variant forms of derivative:
\begin{lyxcode}
\textbackslash usepackage{[}def-file=<filename>{]}\{diffcoeff\}{[}=v4{]}
\end{lyxcode}
This is discussed in �\ref{subsec:The-.def-file}. For the present
document, the call is
\begin{lyxcode}
\textbackslash usepackage{[}def-file=diffcoeff-doc{]}\{diffcoeff\}{[}=v4{]}
\end{lyxcode}
\end{enumerate}

\section{A Rogues' Gallery of derivatives}

\label{sec:Rogues'-gallery}Browsing through texts on statistical
mechanics, relativity and classical mechanics I find the following
choice examples of derivatives `disporting every which way'.

Multi-character variables of differentiation un-parenthesized:

\begin{equation}
\diffp{\frac{\psi}{\Theta}}{\frac{1}{\Theta}},\quad\diffp{E/T}{1/T},\quad\diffp{\ln f}{\ln x_{0}},\quad\diffp\psi{a_{i},\frac{1}{\Theta}},\quad\diffp{\mathcal{L}}{\eta_{,i}^{(r)}}.\label{eq:eg1}
\end{equation}
Multi-character variables of differentiation parenthesized: 
\begin{equation}
\diffp H{\displaystyle \diffp S{q_{k}}[]},\quad\diffp\varepsilon{(1/\Theta)}.\label{eq:eg2}
\end{equation}
Higher-order derivatives where the parentheses do not or sometimes
do include the operator: 
\begin{equation}
\diffp[2]q{\frac{1}{\Theta}},\quad\diffp[2]q{1/\Theta},\quad\diffp[2]\varepsilon{a_{i}},\quad\diff.wrapall.[2]{\phi^{i}(x^{i})}{x^{i}}.\label{eq:eg3}
\end{equation}
Should the $d$ or $\partial$ be included within the parentheses,
as in the last of (\ref{eq:eg3}), or not, as in the others? Logic
says `yes'; practice suggests (generally) `no'.

Indicating a point of evaluation is similarly varied: 
\begin{equation}
\diff.pvrule.\phi\varepsilon[\varepsilon=\varepsilon_{0}],\quad\diff.pvrule.[2]\phi\varepsilon[\varepsilon=\varepsilon_{0}],\quad\diff.psqbra.{b^{\beta}}{a^{\alpha}}[b=0],\quad\diff.paren.uv[v=0].\label{eq:eg4}
\end{equation}
ISO 80000-2 (item 2.11.13) favours the last of these -- parentheses
-- for ordinary derivatives. Presumably, partial derivatives should
follow suit, although parentheses are also used to indicate variables
held constant: 
\begin{equation}
\diffp*{\frac{P}{T}}U[V],\quad\diffp S{N_{2}}[U,V,N_{1}],\quad\diffp S/T[V].\label{eq:eg5}
\end{equation}

Other symbols besides $d$ and $\partial$ are used to denote derivative-like
quantities. From introductory calculus and from classical mechanics
and thermodynamics come $\delta$ and $\Delta$, from fluid mechanics
comes $D$: 
\begin{equation}
\diff.delta.yx,\quad\diff.D.\rho t,\quad\diff.pDelta.UT[V],\quad\diff.Delta.U/T,\quad\diff.delta.{\mathcal{L}}{\eta^{(r)}}.\label{eq:eg6}
\end{equation}

There are those, like the International Organization for Standardization
(ISO), who stipulate (or prefer) an upright `d' for their derivatives:
\begin{equation}
\diff.up.yx.\label{eq:eg11}
\end{equation}
When the differentiand is too big or awkward to sit in the numerator
and is appended to the operator, the $d$ or $\partial$ in the numerator
is generally centred -- but not always. In texts prior to the age
of computerised typesetting one will sometimes find the symbol pushed
to the \emph{left}: 
\begin{equation}
\diff.pleft.*{\diffp{x^{i^{*}}}{x^{k^{*}}}{}}{x^{l^{*}}},\quad\diff.left.*{\left(\frac{m\mathbf{q}_{x}}{\sqrt{1-q^{2}}}\right)}{t}.\label{eq:eg12}
\end{equation}
The observant will note an italic adjustment with the first expression,
so that the $\partial$ in the numerator and the $\partial$ in the
denominator line up in a slanting column, but no such adjustment for
the $d$-s in the second derivative.

Then there is the case when the operator in the numerator differs
from that in the denominator. For instance, in tensor calculus acceleration
is sometimes written 
\[
\diff.nabla.{v^{i}}t=\diff{v^{i}}t+\Gamma_{k\hphantom{i}h}^{\hphantom{k}i}v^{h}\diff{y^{k}}t
\]
where $\nabla v^{i}$ is the `absolute differential' of the velocity
$v^{i}$.

Version 2 or later of the \texttt{diffcoeff} package has the generative
power to cope with all these variations -- see �\ref{sec:Changing-defaults}
-- although it is unlikely an author should need to call on this
capacity to anything like the extent required for this Rogues' Gallery.

Finally and new with version 4 of \texttt{diffcoeff}, is the ability
to insert space before the differentiand, something that had to be
done explicitly by the user before. One way of thinking of a derivative
is as an operator $\diff{}x$ applied to a function $F(x)$ producing
another function $F'(x)$, a \emph{derived }function, the derivative.
Although the original function is included in the numerator of the
differential coefficient, a small space between the $d$ and $F$
feels natural to separate the thing operated on from the thing operating:
\[
F'(x)=\diff!{F(x)}x.
\]
\texttt{diffcoeff} can now produce such spaced derivatives, either
as the default behaviour or at explicit request; see �\ref{subsec:Spacing-before-derivand}.

\section{Ordinary derivatives }

\label{sec:Ordinary-derivatives}Writing\textbf{ }\texttt{\textbackslash diff\{y\}\{x\}}
will produce $\diff{y}{x}$ in text style (i.e., placed between \texttt{\textbackslash (
\textbackslash )} or \texttt{\$ \$}) or 
\[
\diff{y}{x}
\]
in display style (i.e., placed between \texttt{\textbackslash{[}
\textbackslash{]}} ). In fact \texttt{\textbackslash diff yx} (omitting
the braces) will produce these results, with a saving on keystrokes.
The braces are needed only when an argument -- the differentiand
or the variable of differentiation -- is multi-token. 
\begin{itemize}
\item If you want upright `$\mathrm{d}$'s as default, as ISO 80000-2 recommends,
rather than the math-italic `$d$'s used here, this can easily be
done with the package option \texttt{ISO}; see �\ref{sec:Changing-defaults}
on changing default settings. 
\end{itemize}
For inclusion in a line of text you might prefer to use a slash-fraction
form of derivative. That is achieved by inserting a slash, `/', between
numerator and denominator arguments: \texttt{\textbackslash diff\{\textbackslash ln
x\}/x} produces $\diff{\ln x}/x$. Braces are required for the numerator
in this case since it contains more than one token.

\subsection{Spacing before the differentiand}

\label{subsec:Spacing-before-derivand}There are (at least) two different
ways in which we think of derivatives.\footnote{I thank \noun{Hans Sch�lein} for first raising this issue with me
and for subsequent thoughtful comments. } We are all familiar with the argument presented in elementary calculus
books where a curve is shown, and a point on the curve through which
a chord has been drawn. The chord is a side -- the hypotenuse --
of a small right-angled triangle, the other sides having lengths $\delta x$
and $\delta y$ and being parallel to the coordinate axes. The slope
of the chord is $\diffd yx$. By drawing smaller and smaller chords
through the point, the ratio $\diffd yx$ approaches the slope of
the tangent to the curve at the point. We write 
\[
\diff yx
\]
for the limit of $\diffd yx$. It is natural to think of $\dl y$
and $\dl x$ as tiny lengths, like $\delta y$ and $\delta x$, in
which case it would be quite wrong to insert space between the $d$
and the $y$ (let alone the $d$ and the $x$). $dy$ is a single
object, called a differential, and we write expressions like 
\begin{centred}
\verb`\[ dy=\diff yx dx \]` $\Longrightarrow$
\[
\dl y=\diff yx\dl x
\]
\end{centred}
and justly call the `fraction' in this expression a differential
coefficient.

But there is another way of viewing differentiation: as a process
producing (or \emph{deriving}) one function, $f'(x)$, from another,
$f(x)$. Here the sense is of applying $\diff{}x$ to $f(x)$. Although
we include $f(x)$ in the numerator it is not attached to the $d$
and should be separated from it by a small space:
\begin{centred}
\verb`\[ f'(x)=\diff!{f(x)}x \]` $\Longrightarrow$
\[
f'(x)=\diff!{f(x)}x.
\]
\end{centred}
Here the fraction on the right is another name for the derived function
$f'$ and is justly called the derivative of $f$. As you can see
a small space has been inserted between the $d$ and the $f$ in the
numerator. By default the space is \texttt{3 mu} but with the ability
to stretch by \texttt{1 mu} or shrink by\texttt{ 2 mu}{\ttfamily\footnote{In \TeX -speak, \texttt{3mu plus 1mu minus 2mu}. }}
as \TeX{} adjusts lines to fit on the page. (A `mu' is a `math unit'
and is one eighteenth of a quad.) To achieve this result I have inserted
an exclamation mark \verb`!` before the braces delimiting the differentiand. 
\begin{itemize}
\item You may want all or most of your derivatives to have this space and
therefore will not want to be inserting exclamation marks in every
\texttt{\textbackslash diff} command. The \texttt{spaced} package
option switches the default behaviour to spaced derivatives (in which
case the \texttt{!} switch now creates an \emph{un-spaced} derivative).
The size of the space inserted by default can be easily changed; see
�\ref{sec:Changing-defaults}. 
\end{itemize}
The present document uses the un-spaced default. For backwards compatibility
(and perhaps because the author has done this for the last 60 years),
this is the \textquoteleft out-of-the-box\textquoteright{} default
that \texttt{diffcoeff} uses. Authors should make a habit of using
one form predominantly. (There is also the possibility which I haven't
emphasized of spacing multi-token differentiands and leaving single-token
differentiands unspaced; see �\ref{subsec:A-final-flourish}.)

Slash-form derivatives also allow space before the differentiand.
By default this has the same value as applied to the fraction form
of derivative (but can be changed; see �\ref{sec:Changing-defaults}):
\begin{centred}
\verb`$ \diff!{\ln\sin x}/x $` $\Longrightarrow$ $ \diff!{\ln\sin x}/x$.
\end{centred}
The \verb`\diff` command has other optional arguments (e.g. to specify
the order of differentiation) but always the exclamation mark, if
used, immediately precedes the differentiand. It is the positioning
of the differentiand that it affects after all and its placement immediately
before that argument seems natural.

\subsubsection{Ligatures: }

Prior to v.4 of \texttt{diffcoeff}, \texttt{\textbackslash diff fx}
produced the tight pairing evident in 
\[
\frac{df}{dx}.
\]
This was a bug. From version 4.0 of \texttt{diffcoeff}, such ligatures
are prevented: \verb`\[ \diff fx \]` $\Longrightarrow$ \[ \diff fx.\]

\subsubsection{Spacing commands}

\label{subsec:Spacing-commands}The \texttt{diffcoeff} package also
offers four simple spacing commands to fine-tune the display of derivatives
and of other quantities. These are 
\begin{description}
\item [{\texttt{\textbackslash negmu}}] insert spacing of $-1$ mu 
\item [{\texttt{\textbackslash nilmu}}] insert spacing of $0$ mu (cf.
use of an empty brace pair \texttt{\textbf{\{\}}} )
\item [{\texttt{\textbackslash onemu}}] insert spacing of $1$ mu 
\item [{\texttt{\textbackslash twomu}}] insert spacing of $2$ mu 
\end{description}

\subsection{Appending the differentiand: \texttt{\textbackslash diff{*}}}

Some differentiands are too big or awkward to be placed neatly in
the numerator of a derivative and it is natural to append them to
a preceding differential operator. One way to do this is to leave
the numerator argument empty in the \texttt{\textbackslash diff}
command and follow the command with the differentiand. A better way
is to star the \texttt{\textbackslash diff} command. This tells \texttt{diffcoeff}
to append the differentiand. Thus suppose the differentiand is a polynomial,
say $ax^{2}+bx+c$. Add a star (an asterisk) to the \texttt{\textbackslash diff}
command:
\begin{centred}
\verb`\[ \diff*{(ax^2+bx+c)}x \]` $\Longrightarrow$ \[\diff*{(ax^{2}+bx+c)}x\]
\end{centred}
With the \texttt{!} switch or \texttt{spaced} package option, additional
space (by default \texttt{3mu plus 1mu minus 2mu}) is inserted between
the operator and the differentiand:
\begin{centred}
\verb`\[ \diff*!{(ax^2+bx+c)}x \]` $\Longrightarrow$ \[\diff*!{(ax^{2}+bx+c)}x.\]
\end{centred}
A virtue of using an asterisk to append the differentiand is that
if one isn't sure whether a differentiand should be appended or not,
it is an easy matter to simply insert or delete the asterisk to compare
the results. For example, a second derivative is an iterated derivative
-- one in which a derivative forms the differentiand of another derivative:
\begin{centred}
\verb`\[ \diff[2]yx = \diff*{\diff yx}x \]` $\Longrightarrow$ \[ \diff[2]yx=\diff*{\diff yx}x \]
\end{centred}
\noindent which is more elegant to my eye than
\begin{centred}
\noindent \verb`\[ \diff[2]yx = \diff!{\diff yx}x \]` $\Longrightarrow$
\[ \diff[2]yx=\diff!{\diff yx}x \]
\end{centred}
\noindent although whether the \emph{meaning} is clearer is moot.\emph{
}It is easy to switch between the two forms on the right, simply by
inserting or removing the asterisk.

In slash style with the star option, the polynomial example becomes
\begin{centred}
\verb`\[ \diff*{(ax^2+bx+c)}/x \]` $\Longrightarrow$ \[ \diff*{(ax^2+bx+c)}/x \]
\end{centred}
\noindent where the parentheses around the differential operator are
automatically inserted by \texttt{diffcoeff}. Like other elements
of automatic formatting, this is user-adjustable; see �\ref{sec:Changing-defaults}.

With the \texttt{!} switch or \texttt{spaced} package option, this
becomes
\begin{centred}
\verb`\[ \diff*!{(ax^2+bx+c)}/x \]` $\Longrightarrow$ \[ \diff*!{(ax^2+bx+c)}/x \]
\end{centred}

\subsection{Point of evaluation}

\label{subsec:Point-of-evaluation}If you want to specify a point
at which the derivative is evaluated, append a final optional argument.
Note that there \emph{must be no space} before the left square bracket
of the argument:\footnote{In v.1 of \texttt{diffcoeff}, this was brace-delimited. From v.2,
square-brackets have been used and braces deprecated. From v.4, braces
are no longer accepted.}
\begin{centred}
\verb`\[ \diff[2]yx[0] \]` $\Longrightarrow\quad{\displaystyle \diff[2]yx[0]}$
\end{centred}
\noindent If a space does slip in before the final optional argument,
it will not cause a \LaTeX{} error. Instead, the argument will be treated
as a square-bracketed mathematical expression following the derivative,
and typeset as such. 
\begin{itemize}
\item If you prefer to use subscripted \emph{parentheses} around the derivative
to indicate a point of evaluation -- as ISO 80000-2 recommends --
then this can easily be done; see �\ref{sec:Changing-defaults} on
changing default settings. Or, more simply, use the \texttt{ISO} package
option. 
\end{itemize}
Because the slash form spreads the derivative out horizontally, parentheses
are preferred here to indicate a point of evaluation:
\begin{centred}
\verb`$ \diff{\ln sin x}/{sin x}[x=\pi/6] $` $\Longrightarrow$ $ \diff{\ln\sin x}/{\sin x}[x=\pi/6] $.
\end{centred}
\noindent A vertical rule (or `pipe') can easily become too remote
from the opening $d$ of the differential coefficient: $\diff.svrule.{\ln\sin x}/{\sin x}[x=\pi/6]$;
parentheses tie the whole cluster of symbols together.

\subsubsection{Superscripts}

It is easy to add a superscript to a derivative to indicate evaluation
at two points and the difference between the values:
\begin{centred}
\verb`\[ \diff {\sin x}x[0]^{\pi/2} \]` ${\displaystyle \Longrightarrow}$
\[ \diff{\sin x}x[0]^{\pi/2} \]
\end{centred}
\noindent If you want only the superscript, no subscript, include
the final optional argument but leave it empty. Thus, for a particle
of mass $m$ moving along a line, distance $x$ at time $t$, the
kinetic energy is:
\begin{centred}
\noindent \verb`$ \tfrac 12 m \diff x/t[]^2 $` $\Longrightarrow$
$\tfrac{1}{2}m\diff x/t[]^{2}$.
\end{centred}

\subsection{Order of differentiation}

An optional first argument allows the order of differentiation to
be specified. The order need not be a number; an algebraic order of
differentiation is perfectly acceptable as is a mix of the two:
\begin{centred}
\verb`\[ \diff[2]yx \]` $\Longrightarrow$ \[ \diff[2]yx,\]\verb`\[ \diff[n+1]yx \]`
$\Longrightarrow$ \[ \diff[n+1]yx.\]
\end{centred}
As mentioned, the braces can be and have been omitted around the $x$
and $y$; the square brackets around the optional argument, the order
of differentiation, are essential. For a first-order derivative, no
optional argument is needed and entering \texttt{1} as the optional
argument has no effect: 
\begin{centred}
\verb`$ \diff[1]yx $` $\Longrightarrow$$\diff[1]yx$.
\end{centred}
In slash style, \texttt{\$\textbackslash diff{[}2{]}y/x\$} produces
$\diff[2]y/x$, and \texttt{\$\textbackslash diff{[}n+1{]}y/x\$}
produces $\diff[n+1]y/x$.

\subsection{Multi-character variables of differentiation}

Differentiating a function of a function may involve a multi-character
differentiation variable. For instance, to differentiate $\ln\sin x$
in $x$ means forming the product
\begin{centred}
\verb`\[ \diff!{\ln\sin x}{\sin x}\diff{\sin x}x \]` $\Longrightarrow$
\[\diff!{\ln\sin x}{\sin x}\diff{\sin x}x.\]
\end{centred}
\noindent (Although I am mainly using un-spaced differentiands in
this document, a space before the differentiand in the first of these
derivatives -- the \verb`!` switch -- improves the appearance of
the expression to my eye.) 

Forming the \emph{second} derivative of $\ln\sin x$ will now involve
forming (among other quantities)
\begin{centred}
\noindent \verb`\[\diff[2]{\ln\sin x}{\sin x}\]` $\Longrightarrow$
\[\diff[2]{\ln\sin x}{\sin x}\]
\end{centred}
\noindent Parentheses have been inserted automatically by \texttt{diffcoeff}
around $\sin x$ in the denominator to avoid any visual hint that
we are differentiating in the sine of $x^{2}$.

The question is: are the parentheses in the right place? Logically,
no. They should include the $d$: $(d\sin x)^{2}$ -- it is the differential
$d\sin x$ that is of the second order. But as the examples in the
Rogues' Gallery show -- see particularly (\ref{eq:eg3}) -- the
inclination seems to be to do otherwise. This may be because one wants,
in any case, to parenthesise the variable. A second, outer pair of
parentheses then seems too fussy and detracts from comprehending the
symbol `at a glance': 
\[
\diff.wrapall.[2]{f(x)}{(1/x)}.
\]

Customary but illogical notations are familiar in mathematics --
think of the position of the superscripts in an identity like $\sin^{2}\theta+\cos^{2}\theta=1$.
But, like other features of the derivative, the manner of this wrapping
in parentheses of long variables for \emph{higher order} derivatives
is customisable; see �\ref{sec:Changing-defaults}.

For first order derivatives, parenthesising does not occur. If you
want the variable of differentiation to be parenthesised, you need
to insert them yourself:
\begin{centred}
\verb`\[\diff {f(x)}{1/x}, \quad \diff {f(x)}{(1/x)}.\]` $\Longrightarrow$
\[ \diff {f(x)}{1/x}, \quad\diff {f(x)}{(1/x)}.\]
\end{centred}

\section{Partial derivatives}

\noindent \label{sec:Partial-derivatives}Partial derivatives follow
the same pattern as ordinary derivatives, with some extensions. The
command this time is \texttt{\textbackslash diffp}. Thus \texttt{\textbackslash diffp\{F\}\{x\}},
or, with a saving on keystrokes,\texttt{ \textbackslash diffp Fx},
produces $\diffp Fx$ in text style and 
\[
\diffp{F}{x}
\]
in display style. (As for \texttt{\textbackslash diff}, the omission
of braces is possible when the differentiand or the differentiation
variable are single tokens.) As for \texttt{\textbackslash diff},
there is a slash form, generally preferred for inline use, \texttt{\textbackslash diffp
F/x}, displaying as $\diffp F/x$. Given that \texttt{\textbackslash partial}
takes 8 keystrokes to type, the slash form \emph{does }economise on
keystrokes for a partial derivative. 

With either the \verb`spaced` package option (see �\ref{subsec:Package-options})
or the \verb`!` switch a space, defaulting to \verb`3 mu` with some
stretch and shrink, can be inserted before the differentiand: \verb`\[ \diffp!{F(x,y)}x \]`
$\Longrightarrow$ \[ \diffp!{F(x,y)}x. \]

\noindent (But note that if the \verb`spaced` package option is used,
the \verb`!` switch \emph{removes} any extra space.)

Again an optional argument allows the specification of the order of
differentiation which may be numeric or algebraic or a mix of the
two: 
\begin{centred}
\verb`\[ \diffp[3]F/x , \quad \diffp[n]F/x \]` $\Longrightarrow$
\[\diffp[3]F/x , \quad \diffp[n]F/x\] \verb`\[ \diffp[n+1]Fx. \]`
$\Longrightarrow$ \[ \diffp[n+1]Fx. \]
\end{centred}

\subsection{Variables held constant}

In a subject like thermodynamics, there is a need to indicate which
variables are held constant when the differentiation occurs. To show
this, append a final square-bracketed optional argument and ensure
that it follows \emph{immediately} on the preceding mandatory argument.
A space here will detach the argument from the derivative and result
in it being treated as a mathematical expression following the derivative.
Thus to differentiate the entropy $S$ in temperature $T$ while holding
the volume $V$ constant, write
\begin{centred}
\verb`\[ \diffp ST[V] \]` $\Longrightarrow$ \[ \diffp ST[V]. \]
\end{centred}
\noindent In slash form the same expression looks like
\begin{centred}
\noindent \verb`$ \diffp S/T[V] $` $\Longrightarrow$ $ \diffp S/T[V] $.
\end{centred}
This use of a parenthesised, subscripted form to indicate a variable
or variables held constant, leaves open the question: how do we represent
a point of evaluation? ISO 80000-2 makes no recommendation for \emph{partial}
derivatives; presumably we follow the same practice as their recommendation
for ordinary derivatives:
\begin{example}
\verb`\[ \diffp {F(x,y)}x[(0,0)] \]` $\Longrightarrow$ \[ \diffp {F(x,y)}x[(0,0)] \]
\end{example}

However, you may prefer (as many do) to use a vertical rule for this
purpose: 
\[
\diff.pvrule.{F(x,y)}x[(0,0)]
\]
Making this possibility available is discussed in �\ref{sec:Changing-defaults}.

An empty final argument produces a parenthesised derivative with no
subscript,
\begin{centred}
\verb`\[ \diffp yx[] \]` ${\displaystyle \Longrightarrow}$ \[ \diffp yx[] \]
\end{centred}
\noindent which can be useful sometimes. An instance is the writing
of Lagrange's equations of motion in analytic mechanics:
\begin{centred}
\noindent \verb`\[ \diffp L{q_k}-\diff*{\diffp L{\dot{q}_k}[]}t = 0 \]`
$\Longrightarrow$ \[ \diffp L{q_k}-\diff*{\diffp L{\dot{q}_k}[]}t = 0. \]
\end{centred}

\subsubsection{Text-style derivatives}

The \texttt{diffcoeff} package assumes that derivatives formed as
`numerator over denominator' will be used in display-style expressions,
and that the slash form will be used for inline use (text style).
This is the familiar practice in the literature. If one \emph{does}
use the first form in an inline expression where a variable is held
constant, say \texttt{\textbackslash diffp ST{[}V{]}} as here $\diffp ST[V]$,
the result is unsatisfactory, the subscript too tight on the closing
parenthesis and too much space between parentheses and derivative.
The matter is easily resolved using `variant forms' -- see �\ref{subsec:Text-and-script-style}
below -- giving, for our example, $\diff.ptxt.ST[V]$.

\subsection{Appending the differentiand}

\label{subsec:Partial-appending}For a long or awkward differentiand,
it is generally better to \emph{append} it to a preceding differential
operator, rather than create a fractional form with the long expression
in the numerator. As with ordinary derivatives, this is achieved by
adding an asterisk to (i.e. by starring) the \texttt{\textbackslash diffp}
command. 
\begin{centred}
\verb`\[ \diffp*[2]{\Phi(x,y,z)}x \]` $\Longrightarrow$ \[ \diffp*[2]{\Phi(x,y,z)}x. \]
\end{centred}
With the \verb`spaced` package option, or with the \verb`!` switch
if the package option is not used (which is the case for this document),
one gets
\begin{centred}
\verb`\[ \diffp*[2]!{\Phi(x,y,z)}x \]` $\Longrightarrow$ \[ \diffp*[2]!{\Phi(x,y,z)}x \]
\end{centred}
which is an improvement (to my eye).

Alternatively you could leave the first mandatory argument empty and
manually append the differentiand, but by deleting or inserting an
asterisk, it is easy to compare the two forms, differentiand-in-the-numerator,
differentiand-appended, and see which is preferable.

In slash form, parentheses are automatically inserted around the differential
operator when the differentiand is appended,
\begin{centred}
\verb`\[ \diffp*[n]{f(x)}/x \]` $\Longrightarrow$ \[ \diffp*[n]{f(x)}/x, \]
\end{centred}
\noindent although this behaviour can be changed (�\ref{sec:Changing-defaults}
again).

If you wish to both append the differentiand \emph{and} indicate variables
held constant, then the starred form is much the easier way to achieve
this. Thus, to express a relation in thermodynamics,
\begin{centred}
\verb`\[ \diffp*{\frac PT}U[V] = \diffp*{\frac 1T}V[U] \]` $\Longrightarrow$
\[ \diffp*{\frac PT}U[V] = \diffp*{\frac 1T}V[U], \]
\end{centred}
\noindent where the starring automatically takes care of the parentheses
and subscripts.

\subsection{Iterated derivatives}

Derivatives can be the differentiands of derivatives, as seen above
when writing Lagrange's equations,
\begin{centred}
\noindent \verb`\[ \diffp L{q_k}-\diff*{\diffp L{\dot{q}_k}[]}t = 0 \]`
$\Longrightarrow$ \[ \diffp L{q_k}-\diff*{\diffp L{\dot{q}_k}[]}t = 0. \]
\end{centred}
\noindent However, in versions of \texttt{diffcoeff} before version
3.2, an ordinary derivative within a partial derivative rendered as
a partial derivative. That is now corrected, the ordinary derivative
rendering correctly:
\begin{centred}
\verb`\[ \diffp*{\diff{x^\mu}{\lambda}}{x^\sigma} \]` $\Longrightarrow$
\[ \diffp*{\diff{x^\mu}{\lambda}}{x^\sigma}. \]
\end{centred}

\subsection{Mixed partial derivatives}

The new thing with partial derivatives, not present with ordinary
derivatives, is \emph{mixed} partial derivatives, where there is more
than one variable of differentiation. If each variable is differentiated
only to the first order, then it is easy to specify the derivative.
Suppose $F$ is a function of three variables, $x$, $y$ and $z$.
Then
\begin{centred}
\verb`\[ \diffp F{x,y,z} \]` $\Longrightarrow$ \[ \diffp F{x,y,z}, \]

\verb`\[ \diffp !{F(x,y,z)}{x,y,z} \]` $\Longrightarrow$ \[ \diffp !{F(x,y,z)}{x,y,z}. \]
\end{centred}
In the second of these I have inserted space before the differentiand
by using the \verb`!` switch. If the \verb`spaced` package option
was being used, this would have been done automatically.

The variables of differentiation are listed in order in a comma list
forming the second mandatory argument. The total order of differentiation
(3 in this example) is inserted automatically -- \texttt{diffcoeff}
does the calculation. There is also a slash form:
\begin{centred}
\verb`$ \diffp F/{x,y,z} $` $\Longrightarrow$ $ \diffp F/{x,y,z}$.
\end{centred}
If we want to differentiate variables to higher order, then their
orders need to be specified explicitly. To do so use a comma list
for the optional argument:
\begin{centred}
\verb`\[ \diffp[2,3]F{x,y,z} \]` $\Longrightarrow$ \[ \diffp[2,3]F{x,y,z}. \]
\end{centred}
\noindent Notice that the overall order of the derivative -- 6 --
is again automatically calculated and inserted as a superscript on
the $\partial$ symbol in the numerator. In this example, the comma
list of orders has only two members, even though there are three variables.
It is assumed that the orders given in the comma list apply in sequence
to the variables, the first order to the first variable, the second
to the second variable, and so on, and that any subsequent orders
not listed in the optional argument are, by default, 1. Thus we need
to specify only 2 and 3 in the example; the order of differentiation
of $z$ is 1 by default. But you \emph{cannot} use an order specification
like \texttt{{[} , ,2{]}}. Instead write \texttt{{[}1,1,2{]}}.\textbf{
}It is only the \emph{tail} of an order specification which can be
omitted.

\subsubsection{Minutiae of spacing (again)}

\label{subsec:Minutiae-of-spacing}In Chapter 18 of the \emph{The\TeX book},
Knuth suggests inserting a thin space, {\ttfamily\verb`\,`}
(or {\ttfamily\verb`3 mu`}), between differentials in appropriate
contexts, giving as an example $dx\,dy=r\,dr\,d\theta$. In the denominator
of a derivative, however, that degree of extra spacing seems too great
(to my eye), interfering with seeing the derivative `as a whole',
\[
\diff.pwide.F{x,y,z},
\]
especially for the slash-form of derivative: $\diff.pwide.F/{x,y,z}$.
Some spacing is desirable, but less; by default \texttt{diffcoeff}
inserts \verb`2 mu` between the differentials in the fraction form
of derivative and \verb`1 mu`\texttt{ }in the slash form.

Should a differentiation occur to higher order and so a variable acquire
a superscript, an adjustment is made to the extra spacing. By default
\verb`1 mu` is subtracted from the default spacing. Thus in 
\[
\diffp[2]F{x,y,z},
\]
\verb`2 mu` of spacing is inserted between the $\partial y$ and
$\partial z$, but because the superscript already provides some separation
between them, only \verb`1 mu` is inserted between $\partial x^{2}$
and $\partial y$.

The values used for the spacing and its adjustment in the presence
of a superscript can be changed by the user; see �\ref{sec:Changing-defaults}.
Some other minutiae of spacing are discussed below in �\ref{subsec:Comma-list-of-vars}
when the variables themselves are super- or subscripted (as happens
in tensor calculus, for instance).

Note that the \verb`!` switch, if used, is inserted \emph{immediately
before} the differentiand. It is the placement of the differentiand,
after all, that it is affecting:
\begin{centred}
\verb`\[ \diffp[2]!F{x,y,z}  \]` $\Longrightarrow$ \[ \diffp[2]!F{x,y,z}. \]
\end{centred}

\subsubsection{Algebraic orders of differentiation}

Calculation of the overall order of differentiation still occurs even
when some or all of the individual orders are algebraic variables
rather than numbers. For example, differentiating in three variables
to orders \texttt{1, m+1, m-1},
\begin{centred}
\verb`\[ \diffp[1,m+1,m-1]{F(x,y,z)}{x,y,z} \]` $\Longrightarrow$
\[ \diffp[1,m+1,m-1]{F(x,y,z)}{x,y,z}. \]
\end{centred}

\subsubsection{Error messages}

Should you specify \emph{more} orders in the comma list of the order
specification than there are variables, compilation will fail and
an error message will be sent to the terminal and \LaTeX{} log . For
example, if on (say) line 53 of my document I specify \texttt{{[}1,m-1,m+1,2{]}}
for the orders of differentiation but list only \texttt{\{x,y,z\}}
for the variables, the message will be 
\begin{lyxcode}
!~Package~diffcoeff~Error:~4~orders~of~differentiation~

specified~for~3~variables;~orders~{[}1,m+1,m-1,2{]}

(on~line~53)~for~variables~x,y,z.
\end{lyxcode}
Should you try using a \emph{number} raised to a power in an order
specification, it will cause an error. For example, if on line 53
(say) of my document I try to form \texttt{\textbackslash diffp{[}2\textasciicircum 2{]}f\{x,y,z\}}
compilation will fail with the message 
\begin{lyxcode}
!~Package~diffcoeff~Error:~number~followed~by~\textasciicircum ~in~the~

order~specification~{[}2\textasciicircum n,1,1{]}~on~line~53.~Diffcoeff~

cannot~calculate~the~overall~order~of~differentiation~

in~this~case.~Use~the~order-override~option~to~enter~

the~overall~order.
\end{lyxcode}
The order-overide option mentioned here is discussed next. Similarly,
you cannot use multiplication (\texttt{\textbackslash times} or \texttt{{*}})\texttt{
}or division (\texttt{/} or \texttt{\textbackslash div}) signs after
\emph{numbers} in an order specification; only \texttt{+} or \texttt{-}
or a left parenthesis (see below) are allowed.

\subsubsection{The order-override option}

\noindent Perhaps the differentiations are to orders \texttt{{[}2km,m-1,m+1{]}}:
\begin{centred}
\noindent \verb`\[ \diffp[2km,m-1,m+1]{F(x,y,z)}{x,y,z} \]` $\Longrightarrow$
\[ \diffp[2km,m-1,m+1]{F(x,y,z)}{x,y,z}. \]
\end{centred}
\noindent Here the overall order is presented as \texttt{2km+2m}.
You might prefer this to be presented as \texttt{2m(k+1)}. Although
\texttt{diffcoeff} takes some steps to present the overall order appropriately,
it is not a computer algebra system and does not factorise expressions.
If you want to present the order in a manner distinct from that presented
by \texttt{diffcoeff}, use the \emph{order-override option}.

This is a second optional argument immediately following the order
specification. For our last example, filling the override option with
\texttt{2m(k+1)} gives the desired result:
\begin{centred}
\verb`\[ \diffp[2km,m-1,m+1][2m(k+1)]{F(x,y,z)}{x,y,z} \]` $\Longrightarrow$
\[ \diffp[2km,m-1,m+1][2m(k+1)]{F(x,y,z)}{x,y,z}. \]
\end{centred}
\noindent As another example, left to its own devices, \texttt{diffcoeff}
produces
\begin{centred}
\noindent \verb`\[ \diffp[m/2+n/2,m/2,n/2]F{x,y,z} \]` $\Longrightarrow$
\[ \diffp[m/2+n/2,m/2,n/2]F{x,y,z}, \]
\end{centred}
\noindent whereas we would like the total order to be presented as
$m+n$. Using the override option,
\begin{centred}
\noindent \verb`\[ \diffp[m/2+n/2,m/2,n/2][m+n]F{x,y,z} \]` $\Longrightarrow$
\[ \diffp[m/2+n/2,m/2,n/2][m+n]F{x,y,z}. \]
\end{centred}
The order-override option does exactly that: overrides the presentation
of the calculated order with the manually given one. In fact the calculation
algorithm does not get called at all when the override option is used
so that one can in this way present the total order in whatever manner
one wishes or, indeed, add essentially arbitrary material as a superscript
to the $\partial$ symbol in the numerator.

\subsubsection{Comma list of variables of differentiation}

\label{subsec:Comma-list-of-vars}In tensor calculus the differentiations
are almost always in terms of super- or subscripted coordinates. In
many other contexts this is the case too -- the reciprocal of the
temperature in thermodynamics or generalized coordinates in analytical
mechanics. This is why a comma list is used in \texttt{diffcoeff}
for specifying variables of differentiation for mixed partial derivatives.
Although it would be nice to write the minimal \texttt{\{xy\}} rather
than \texttt{\{x,y}\} when two variables $x$ and $y$ are involved,
the extra writing is trivial and the comma list allows a simpler handling
of multi-character variables. For instance in tensor calculus we get
expressions like
\begin{centred}
\verb`\[ \diffp{A_i}{ x^j,x^k } \]` $\Longrightarrow$ \[ \diffp{A_i}{ x^j,x^k }. \]
\end{centred}
\noindent It is easier to write \texttt{\{x\textasciicircum j,x\textasciicircum k\}}
here than, say, \texttt{\{\{x\textasciicircum j\}\{x\textasciicircum k\}\}}
to distinguish the variables. It does mean that should the variable
of differentiation include a comma then that comma needs to be enclosed
in braces. There are plenty of instances of this out in the world
(see, e.g., the last equation of (\ref{eq:eg1})) but it is overall
a rare occurrence.

\paragraph*{Minutiae of spacing (yet again):}

In �\ref{subsec:Minutiae-of-spacing} above, I discussed a slight
reduction in the space inserted between the terms occurring in the
denominator of a mixed partial derivative when a higher order differentiation
is involved. The superscript appearing on a differentiation variable
in that case \emph{of itself} introduced a spacing adjustment between
the terms. But the present discussion involves only first order differentiations
and no such reduction is automatically made by \texttt{diffcoeff}.
However it is still possible to explicitly make such an adjustment
with the \texttt{\textbackslash negmu} command introduced earlier
(�\ref{subsec:Spacing-commands}), which inserts \texttt{-1 mu} of
spacing. For our example, in
\begin{centred}
\verb`\[ \diffp{A_i}{ x^j\negmu,x^k } \]` $\Longrightarrow$ \[ \diffp{A_i}{ x^j\negmu,x^k } \]
\end{centred}
\noindent the \texttt{\textbackslash negmu} decreases the spacing
between the terms from the default \verb`2 mu` to \verb`1 mu`.

\subsubsection{Overkill territory}

Two previous examples illustrate limitations of the algorithm that
calculates the overall order of differentiation: \texttt{2m/2+2n/2}
is not simplified to \texttt{m+n} and \texttt{2km+2m} is not factorised
to \texttt{2m(k+1)}. But there is much that the algorithm \emph{can}
handle -- for instance, the simple use of parentheses:
\begin{centred}
\verb`\[ \diffp[2m-(k+1),2(k+1)-m]{F(x,y,z)}{x,y,z} \]` $\Longrightarrow$
\[ \diffp[2m-(k+1),2(k+1)-m]{F(x,y,z)}{x,y,z}. \]
\end{centred}

\paragraph*{Dynamic use of parentheses}

\noindent For parenthetic expressions to be evaluated as in this example
-- the \emph{dynamic} use of parentheses -- the left parenthesis
must be preceded at most by a sign or a number; the right parenthesis
must be followed at most by a sign.

If a right parenthesis is followed by a \emph{variable}, say by \texttt{m}
as in the order spec. \texttt{{[}(2n+1)m,(2n-1)m{]}}, it will throw
an error and halt compilation. A message will be sent to the terminal
and the \LaTeX{} log like the following (which assumes the order spec.
was on line 53 of the document): 
\begin{lyxcode}
!~Package~diffcoeff~Error:~)~followed~by~m~in~the~

order~specification~{[}(2n+1)m,(2n-1)m{]}~on~line~53.~

Diffcoeff~cannot~calculate~the~overall~order~of~

differentiation~in~this~case.~Use~the~order-override~

option~to~enter~the~overall~order.
\end{lyxcode}
This is a limitation on the dynamic use of parentheses -- but they
\emph{can} be nested.

\paragraph*{Static use of parentheses}

If a left parenthesis is preceded by a \emph{variable} (i.e., not
a sign or a number) this signals to \texttt{diffcoeff} the \emph{static}
use of parentheses, something like $f(k)$, `function of $k$'.
No attempt is made to evaluate what is between them and they are treated
simply as an extension of the variable. For example,
\begin{centred}
\verb`\[ \diffp[f(k+1)+1,f(k-1)-1]{F(x,y)}{x,y} \]` $\Longrightarrow$
\[ \diffp[f(k+1)+1,f(k-1)-1]{F(x,y)}{x,y}. \]
\end{centred}
\noindent In the static case you \emph{can} follow the right parenthesis
by a variable without generating an error.\emph{ }You can nest them,
and you can include static parentheses within a dynamic pair; for
example,
\begin{centred}
\noindent \verb`\[ \diffp[2(3+f(k))+1,1-3(f(k)-2)]{F(x,y)}{x,y} \]`
$\Longrightarrow$ \[ \diffp[2(3+f(k))+1,1-3(f(k)-2)]{F(x,y)}{x,y}. \]
\end{centred}
\noindent However, the reverse is not possible: you can't have dynamic
parentheses within a static pair.

\paragraph*{Other refinements}

Exponents and subscripts on a \emph{variable} are fine in an order
specification, so long as the exponent or subscript consists of a
\emph{single} token:
\begin{centred}
\verb`\[ \diffp[m^2+2(k-1),m^2-(k+1)]F{x,y,z,w} \]` $\Longrightarrow$
\[ \diffp[m^2+2(k-1),m^2-(k+1)]F{x,y,z,w}. \]
\end{centred}
\noindent Braced arguments containing \emph{multiple} tokens as exponents
or subscripts to variables will generally not halt compilation but
will usually give nonsensical results, as will \emph{signs} treated
as superscripts or subscripts. Neither circumstance is checked for
by \texttt{diffcoeff}.

\paragraph*{Override}

There are limitations on what order specifications the \texttt{diffcoeff}
package can `digest'; equally, it can digest a wide variety of such
constructs, but it is \emph{not} a computer algebra system. In all
those cases where it fails to calculate or present a correct total
order, the order-override option is available. Yes, this is not as
convenient as having the overall order calculated automatically but
(let's remind ourselves) we are deep in overkill territory. Mixed
partial derivatives are used far less often than the pure derivatives,
and when they \emph{are} used it is nearly always to orders 1 or 2
in the variables. Mixed partial derivatives to exotic orders of differentiation
are rarely used, so that the limitations of the calculational algorithm
are of little real moment -- and the override option is always available
for such cases.

\subsection{Parentheses around multi-character variables}

In thermodynamics and statistical mechanics one may want to differentiate
in the reciprocal of the temperature, $1/T$ (or $1/\Theta$):
\begin{centred}
\verb`\[ \diffp[2]q{\frac 1\Theta} \]` $\Longrightarrow$ \[ \diffp[2]q{\frac 1\Theta}. \]
\end{centred}
\noindent In this case and for other \emph{higher order} derivatives
of multi-character variables of differentiation, the parentheses are
inserted automatically by \texttt{diffcoeff}. Precisely where parentheses
should be placed is moot. The placement in this example is not strictly
logical, although it feels intuitive, but the placement can be customised
(�\ref{sec:Changing-defaults}).

Parentheses are automatically inserted like this only for higher order
derivatives. When the differentiation is to first order, parenthesising
is up to the user:
\begin{centred}
\verb`\[ \diffp q{(\frac 1\Theta),V} \]` $\Longrightarrow$ \[ \diffp q{(\frac 1\Theta),V}. \]
\end{centred}

\subsection{Jacobians}

\texttt{diffcoeff} provides a command \texttt{\textbackslash jacob}
for constructing Jacobians. For example
\begin{centred}
\verb`\[ \jacob{u,v,w}{x,y,z} \]` $\Longrightarrow$ \[ \jacob{u,v,w}{x,y,z}. \]
\end{centred}
The comma lists can contain any number of variables. \texttt{\textbackslash jacob}
does \emph{not} check that the two arguments contain the same number
of variables, so it is perfectly possible to form an object like \texttt{\textbackslash jacob\{u,v,w\}\{x,y\}}
which as far as I know has no meaning.

\section{Changing defaults; variant forms}

\label{sec:Changing-defaults}To write the range of different examples
displayed in the Rogues' Gallery (�\ref{sec:Rogues'-gallery}) I have
had to make extensive use of forms of derivative other than the default.
\texttt{diffcoeff} is built on the facilities offered by the \texttt{xtemplate}
package (included in the \LaTeX 3 bundle \texttt{l3packages)}. These
facilities are harnessed by means of a command, \texttt{\textbackslash diffdef},
and a further optional argument of the \texttt{\textbackslash diff}
command.\texttt{ }

\subsection{Default values: ordinary derivatives}

Table~\ref{tab:Ordinary-derivatives} lists the keys available for
forming derivatives and the default values\footnote{Note that a mu is a `math unit', 1/18 of a quad.}
they have been assigned. These default values have been chosen to
coincide with those relevant for \emph{ordinary} derivatives -- apart
from the keys \texttt{denom-term-sep}, \texttt{/-denom-term-sep},
\texttt{term-sep-adjust} and \texttt{/-term-sep-adjust} which are
ignored for ordinary derivatives but apply to the case of mixed partial
derivatives when there is more than one variable of differentiation.
Keys with an opening slash, /, apply only to the slash form of the
derivative; keys with an opening asterisk, {*}, apply only when the
differentiand is appended.

Note that these settings are, in general, font dependent. The values
given are (in the author's opinion) appropriate for the default \LaTeX{}
math fonts. There are also likely to be variations required for whether
a derivative is used in a display-style or text-style or script-style
expression. That matter is discussed below in �\ref{subsec:Text-and-script-style}.
All values specifying a space require the unit (\texttt{mu}) to be
given; a number alone does not suffice.
\noindent \begin{center}
\begin{table}
\centering
\centering{}\caption{{\small Defaults (ordinary derivatives})\protect\label{tab:Ordinary-derivatives}}
\begin{tabular}{lr}
\toprule 
{\small key} & {\small default}\tabularnewline
\midrule
{\small op-symbol} & {\small\texttt{d}}\tabularnewline
{\small op-symbol-alt} & {\small\texttt{op-symbol}}\tabularnewline
{\small op-order-sep} & {\small\texttt{1 mu}}\tabularnewline
{\small derivand-sep} & {\small\texttt{3 mu plus 1 mu minus 2 mu}}\tabularnewline
{\small /-derivand-sep} & {\small\texttt{derivand-sep}}\tabularnewline
{\small{*}-derivand-sep} & {\small\texttt{derivand-sep}}\tabularnewline
{\small{*}/-derivand-sep} & {\small\texttt{/-derivand-sep}}\tabularnewline
{\small denom-term-sep} & {\small\texttt{2 mu}}\tabularnewline
{\small /-denom-term-sep} & {\small\texttt{1 mu}}\tabularnewline
{\small term-sep-adjust} & {\small\texttt{-1 mu}}\tabularnewline
{\small left-delim} & {\small\texttt{\textbackslash left .}}\tabularnewline
{\small right-delim} & {\small\texttt{\textbackslash right |}}\tabularnewline
{\small /-left-delim} & {\small\texttt{(}}\tabularnewline
{\small /-right-delim} & {\small\texttt{)}}\tabularnewline
{\small elbowroom} & {\small\texttt{0 mu}}\tabularnewline
{\small /-elbowroom} & {\small\texttt{0 mu}}\tabularnewline
{\small subscr-nudge} & {\small\texttt{0 mu}}\tabularnewline
{\small /-subscr-nudge} & {\small\texttt{0 mu}}\tabularnewline
{\small long-var-wrap} & {\small\texttt{d(v)}}\tabularnewline
{\small{*}/-op-wrap} & {\small\texttt{true}}\tabularnewline
{\small{*}-op-left} & {\small\texttt{false}}\tabularnewline
{\small{*}-italic-nudge} & {\small\texttt{3 mu}}\tabularnewline
\bottomrule
\end{tabular}
\end{table}
\par\end{center}
\begin{description}
\item [{op-symbol}] the operator symbol; for ordinary derivatives, generally
one of \texttt{d} or \texttt{\textbackslash mathrm\{d\}}, \texttt{D}
or \texttt{\textbackslash mathrm\{D\}}, \texttt{\textbackslash delta}
or \texttt{\textbackslash Delta}; for partial derivatives \texttt{\textbackslash partial};
default = \texttt{d}
\item [{op-symbol-alt}] if different from \textbf{op-symbol} then used
in the denominator while \textbf{op-symbol} is used in the numerator;
e.g. for $\diff.nabla.{v^{i}}t$, \texttt{op-symbol = \textbackslash nabla}
and \texttt{op-symbol-alt = d}; otherwise (and usually) defaults to
\textbf{op-symbol}
\item [{op-order-sep}] extra horizontal space added between the op-symbol
and the superscripted order of differentiation in higher order derivatives;
for the math-italic forms compare $d^{2}$ with $d^{\mkern1mu 2}$,
$\partial^{2}$ with $\partial^{\mkern1mu 2}$ where the first symbol
in each case has no extra space and the second has an extra 1 mu;
default = \texttt{1 mu}
\item [{derivand-sep}] horizontal space added before the differentiand
(derivand) if the \verb`spaced` package option is used, or by the
\verb`!` switch if it is not; the default has some stretch and shrink;
default = \verb`3mu plus 1mu minus 2mu`
\item [{/-derivand-sep}] for the slash form of derivative, horizontal space
added before the differentiand (derivand) if the \verb`spaced` package
option is used, or by the \verb`!` switch if it is not; default =
\verb`derivand-sep`
\item [{{*}-derivand-sep}] when the derivand is appended, horizontal space
added before the differentiand (derivand) if the \verb`spaced` package
option is used, or by the \verb`!` switch if it is not; default =
\verb`derivand-sep`
\item [{{*}/-derivand-sep}] for the slash form of derivative when the derivand
is appended, horizontal space added before the differentiand (derivand)
if the \verb`spaced` package option is used, or by the \verb`!`
switch if it is not; default = \verb`/-derivand-sep`
\item [{denom-term-sep}] (mixed partial derivatives only) horizontal spacing
inserted between the differentials in the denominator of a mixed partial
derivative to avoid a solid cluster like $\partial x\partial y\partial z$;
with the default 2 mu this is $\dl.p.x\dl.p.2y\dl.p.2z$; default
= \texttt{2 mu}
\item [{/-denom-term-sep}] (mixed partial derivatives only) horizontal
spacing inserted between the differentials in the denominator of a
slash-form mixed partial derivative; because a slash-form derivative
is already spread out horizontally, the default spacing is less than
for the \texttt{\textbackslash frac} form derivative; default = \texttt{1
mu}
\item [{term-sep-adjust}] (mixed partial derivatives only) adjustment (i.e.
reduction) to \textbf{denom-term-sep} or \textbf{/-denom-term-sep}
when differentiation in a variable occurs to an order other than 1;
if, e.g., $\dl.p.x^{2}\dl.p.1y\dl.p.2z$ is the denominator of a mixed
partial derivative in three variables, because of the superscript
the spacing between $\partial x^{2}$ and $\partial y$ is reduced
by\textbf{ term-sep-adjust} from the spacing between $\partial y$
and $\partial z$; default = \texttt{-1 mu}
\item [{left-delim}] the left member of a delimiter pair wrapping the derivative,
the right member of which is subscripted to indicate a point of evaluation
or variables held constant; default = \texttt{\textbackslash left
.}
\item [{right-delim}] the right member of a delimiter pair wrapping the
derivative and subscripted to indicate a point of evaluation or variables
held constant; default = \texttt{\textbackslash right |}
\item [{/-left-delim}] for the slash form of derivative, the left member
of a delimiter pair wrapping the derivative and subscripted to indicate
a point of evaluation or variables held constant; default = \texttt{(}
\item [{/-right-delim}] for the slash form of derivative, the right member
of a delimiter pair wrapping the derivative, the right member of which
is subscripted to indicate a point of evaluation or variables held
constant; default = \texttt{)}
\item [{elbowroom}] adjustment to the whitespace between the left and right
delimiters and the enclosed derivative; negative values reduce the
space; default = \texttt{0 mu}
\item [{/-elbowroom}] adjustment to the whitespace between the left and
right delimiters and the enclosed slash-form derivative; default =
\texttt{0 mu} 
\item [{subscr-nudge}] horizontal adjustment of the subscript's placing
relative to the \textbf{right-delim}iter, e.g., a negative value compensates
for the curving inwards of a large right parenthesis; may be font
dependent; default = \texttt{0 mu}
\item [{/-subscr-nudge}] for the slash form of derivative, horizontal adjustment
of the subscript's placing relative to the /-\textbf{right-delim}iter;
may be font dependent; default = \texttt{0 mu} 
\item [{long-var-wrap}] to avoid ambiguity in higher order derivatives
it may be advisable to wrap multi-token variables of differentiation
in parentheses; default = \texttt{d(v)}; the choices are 
\begin{description}
\item [{\texttt{dv}}] no wrapping, e.g. $dx_{i}^{2}$ or $d\frac{1}{\Theta}^{2}$,
$\partial x_{i}^{2}$ or $\partial\frac{1}{\Theta}^{2}$, 
\item [{\texttt{d(v)}}] wrap the variable only, e.g. $d(x_{i})^{2}$ or
$d(\frac{1}{\Theta})^{2}$, $\partial(x_{i})^{2}$ or $\partial(\frac{1}{\Theta})^{2}$
\item [{\texttt{(dv)}}] wrap the op-symbol and variable, e.g. $(dx_{i})^{2}$
or $(d\frac{1}{\Theta})^{2}$, $(\partial x_{i})^{2}$ or $(\partial\frac{1}{\Theta})^{2}$
\end{description}
\item [{{*}/-op-wrap}] a choice of \texttt{true} or \texttt{false} for
slash forms of the derivative when the differentiand is appended,
dictating whether the derivative is wrapped in parentheses, as here
$\diffp*{F(x,y)}/x$, or not; default = \texttt{true}
\item [{{*}-op-left}] a choice of \texttt{true} or \texttt{false} indicating
whether the op-symbol is left-aligned or not when the differentiand
is appended; generally it is centred; does not apply to slash forms
of the derivative; default = \texttt{false}
\item [{{*}-italic-nudge}] if \textbf{{*}-op-left} is \texttt{true}, makes
an italic adjustment in the numerator, so that the op-symbols in numerator
and denominator align in the same slanting column; for an upright
\texttt{d} this would be set to \texttt{0 mu}; default = \texttt{3
mu}
\end{description}

\subsection{ISO defaults}

\label{subsec:ISO-defaults}\begin{wraptable}[10]{o}{0.4\columnwidth}%
\centering{}\vspace{-5.35ex}
 \caption{{\small ISO default changes}\protect\label{tab:ISO-setting-changes}}
\abovetopsep=.5ex %
\begin{tabular}{lr}
\toprule 
{\small key} & {\small default}\tabularnewline
\midrule 
{\small op-symbol} & {\small\texttt{\textbackslash mathrm\{d\}}}\tabularnewline
{\small op-order-sep} & {\small\texttt{0 mu}}\tabularnewline
{\small left-delim} & {\small\texttt{\textbackslash left (}}\tabularnewline
{\small right-delim} & {\small\texttt{\textbackslash right )}}\tabularnewline
{\small subscr-nudge} & {\small\texttt{-6 mu}}\tabularnewline
\bottomrule
\end{tabular}\end{wraptable}%
You may not like the default settings that come with \texttt{diffcoeff}.
The package does not follow ISO 80000-2 -- it does not use upright
`d's nor does it wrap an ordinary differential coefficient in subscripted
parentheses to indicate a point of evaluation. Both `defects' can
be remedied by calling the package with the option \texttt{ISO}:\footnote{One can also use \texttt{ISO=true} to turn the option on and \texttt{ISO=false
}to turn the option off. } 
\begin{lyxcode}
\textbackslash usepackage{[}ISO{]}\{diffcoeff\}{[}=v4{]}
\end{lyxcode}
The uppercase is essential -- an option \texttt{iso} is not recognised.
The \texttt{ISO} option results in changes to the default settings
of Table~\ref{tab:Ordinary-derivatives} as listed in Table~\ref{tab:ISO-setting-changes}.
Any settings not mentioned in Table~\ref{tab:ISO-setting-changes}
retain the values presented in Table~\ref{tab:Ordinary-derivatives}.
Note that the subscript nudge figure specified here is \emph{not}
part of the standard, which makes no recommendation about the subscript's
positioning. But: the \texttt{-6 mu} figure with a default or latin
modern font gives a better representation of what is displayed in
the standard than a zero figure.

Because the `d' is upright with the \texttt{ISO} option, no extra
space is required between the symbol and the superscript in a higher
order derivative. Hence the zero value for the \texttt{op-order-sep}.
ISO recommends subscripted parentheses to indicate a point of evaluation.
Hence the other entries in the table. Because a large right parenthesis
(penultimate setting) bends inwards, to the left, a negative value
for the last entry ensures the subscript does not become detached
from the derivative, looking lost in a sea of whitespace.

Note that the \texttt{ISO} option will also produce upright `D's in
derivatives formed from `D'; see �\ref{subsec:D-delta-Delta} below.

\subsection{Partial derivatives}

\begin{wraptable}{o}{0.4\columnwidth}%
\centering{}\vspace{-5.35ex}
 \caption{{\small Default changes for partial derivatives}\protect\label{tab:Partial-deriv-defaults}}
\abovetopsep=.5ex %
\begin{tabular}{lr}
\toprule 
{\small key} & {\small default}\tabularnewline
\midrule 
{\small op-symbol} & {\small\texttt{\textbackslash partial}}\tabularnewline
{\small left-delim} & {\small\texttt{\textbackslash left (}}\tabularnewline
{\small right-delim} & {\small\texttt{\textbackslash right )}}\tabularnewline
{\small subscr-nudge} & {\small\texttt{-6 mu}}\tabularnewline
\bottomrule
\end{tabular}\end{wraptable}%
The default values given in Table~\ref{tab:Ordinary-derivatives},
when they are relevant, apply to \emph{ordinary} derivatives. For
partial derivatives, some defaults change. These are listed in Table~\ref{tab:Partial-deriv-defaults}.
All other keys take the default values of Table~\ref{tab:Ordinary-derivatives}. 

The last three entries here reflect the use of subscripted parentheses
with partial derivatives to indicate variables held constant, for
instance in the Maxwell relations of thermodynamics, one of which
is 
\[
\diffp SV[T]=\diffp PT[V]
\]


\subsection{Setting your own defaults: \texttt{\textbackslash diffdef}}

\label{subsec:diffdef}Versions 2 and later of the \texttt{diffcoeff}
package provide a command, \texttt{\textbackslash diffdef}, that
enables users to set their own defaults.\texttt{ }For example, if
you wish to use upright `d's but not follow ISO's use of subscripted
parentheses to indicate a point of evaluation, enter in the preamble
of your document the command
\begin{lyxcode}
\textbackslash diffdef~\{\}~~~~

~~\{~~~~~~

~~~~op-symbol~~~~=~\textbackslash mathrm\{d\},~~~~

~~~~op-order-sep~=~0~mu

~~\}
\end{lyxcode}
\noindent Since a list of settings, like this one, is a comma-\emph{separated}
list, no comma is required for the last entry. That entry is a consequence
of the first: upright symbols do not require any extra separation
between the `d' and the superscript in a higher order derivative.

The other point to note is the empty pair of braces after the \texttt{\textbackslash diffdef}
command. \emph{They matter}. Their emptiness is what determines that
it is the \emph{default} values that are changed. If they contain
some content, then that content provides a \emph{name} for the particular
set of values in the following list. The \texttt{diffcoeff} package
exploits this facility to cope with the wide variety of forms displayed
in the Rogues' Gallery of �\ref{sec:Rogues'-gallery}.

\subsubsection{Space before the differentiand}

\label{subsec:A-final-flourish}\begin{wraptable}{o}{0.6\columnwidth}%
\begin{centering}
\caption{Keys for spacing the derivand}
\medskip{}
\begin{tabular}{lr}
\toprule 
{\small key} & {\small default}\tabularnewline
\midrule
{\small derivand-sep} & {\small\texttt{3mu plus 1mu minus 2mu}}\tabularnewline
{\small{*}-derivand-sep} & {\small\texttt{derivand-sep}}\tabularnewline
{\small /-derivand-sep} & {\small\texttt{derivand-sep}}\tabularnewline
{\small{*}/-derivand-sep} & {\small\texttt{ /-derivand-sep}}\tabularnewline
\bottomrule
\end{tabular}
\par\end{centering}
\end{wraptable}%
The insertion of a small space before the differentiand is effected
by the \verb`!` key inserted immediately before the differentiand
argument in the \verb`\diff` command. You may wish to make the insertion
of this space the \emph{default} behaviour. This is done by using
the \verb`spaced` package option (which makes the \verb`!` switch
now reverse this new default and put \emph{no} extra space before
the derivand). The amount of space inserted is, by default, \verb`3mu plus 1mu minus 2mu`,
meaning the space is generally $3$mu but can stretch to $4$mu or
shrink to $1$mu as \TeX{} strives to fit content in a line or on the
page.

Perhaps this doesn't suit. You may want a fixed space here, with no
stretch or shrink. The key to change is \verb`derivand-sep`. By default,
this setting applies not only to the fraction form of derivative,
but also to the slash form and to when the derivand is appended. 

If you feel a little less space should be used for slash derivatives,
then the key to change is \verb`/-deriv-sep`. This changed value
will also be used for an appended derivand in a slash derivative. 

Thus to meet both wishes you might put in the preamble of your document
something like
\begin{lyxcode}
\textbackslash diffdef~\{\}~~~~

~~\{~~~~~~

~~~~derivand-sep~~~=~3~mu,~~~~

~~~~/-derivand-sep~=~2~mu

~~\}
\end{lyxcode}
This will insert a fixed space of $3$mu before the differentiand
in both the numerator and when appended in the fraction form of derivative,
and a fixed space of $2$mu in the slash form of derivative, both
in the numerator and when appended.

If you want a different spacing when the derivand is appended, the
keys to change are \verb`*-derivand-sep` and \verb`*/-derivand-sep`.

\paragraph{Selective spacing}

I have treated the \verb`spaced` package option thus far as if it
were an \verb`ON/OFF` switch and, indeed, the presence of the package
option \verb`spaced` behaves as \verb`ON` and its absence as \verb`OFF`.
Internally, however, \verb`spaced` is equivalent to \verb`spaced=1`
and its absence to \verb`spaced=0`. Entering \verb`spaced=n` in
the package option where \verb`n` is a positive integer is equivalent
to entering \verb`spaced=1` (and hence to simply entering \verb`spaced`),
but if \verb`n` is a negative integer, a new effect is produced. 

Entering \verb`spaced=-1` (or any negative integer) as a package
option will insert a space (by default \verb`3mu plus 1mu minus 2mu`)
before the differentiand provided the differentiand\emph{ is longer
than a single token} but will insert no space before single-token
differentiands. The switch \verb`!` reverses this behaviour.

\subsection{Variant forms}

For this package I needed a number of variant forms to illustrate
the wide variety of ways in which derivatives are displayed. The \texttt{\textbackslash diffdef}
command in which the first argument is \emph{filled} provides one
half of the means of doing this. I've given the single-letter name
\texttt{p} to the following settings:
\begin{lyxcode}
\textbackslash diffdef~\{~p~\}

~~\{

~~~~op-symbol~~~~=~\textbackslash partial~,

~~~~left-delim~~~=~\textbackslash left~(~~,

~~~~right-delim~~=~\textbackslash right~)~,

~~~~subscr-nudge~=~-6~mu

~~\}
\end{lyxcode}
The second half of providing variant forms is to insert this name,
\texttt{p}, between dots (periods, full stops) as the \emph{first}
argument of the \texttt{\textbackslash diff} command. Thus, repeating
an example at the end of �\ref{subsec:Partial-appending},
\begin{centred}
\verb`\[ \diff.p.*{\frac PT}U[V] = \diff.p.*{\frac 1T}V[U] \]` $\Longrightarrow$
\[ \diff.p.*{\frac PT}U[V] = \diff.p.*{\frac 1T}V[U] \]
\end{centred}
\noindent The effect is exactly the same as previously, when the \texttt{\textbackslash diffp}
command was used. Indeed, \texttt{diffcoeff} identifies \texttt{\textbackslash diffp}
with \texttt{\textbackslash diff.p.}:
\begin{lyxcode}
\textbackslash NewDocumentCommand~\textbackslash diffp~\{~\}~\{~\textbackslash diff.p.~\}
\end{lyxcode}
Note that this identification of \texttt{\textbackslash diffp} with
\texttt{\textbackslash diff.p.} means there is no equivalent dot-delimited
argument available for \texttt{\textbackslash diffp}. The \emph{dot-delimited
argument applies only to} \texttt{\textbackslash diff}.

For example, to illustrate the upright-d form of derivative, without
changing the default math-italic form (which I prefer), one might
enter in the preamble 
\begin{lyxcode}
\textbackslash diffdef~\{~up~\}~

~~\{

~~~~op-symbol~~~~=~\textbackslash mathrm\{d\},

~~~~op-order-sep~=~0~mu

~~\}
\end{lyxcode}
Apart from the \emph{key = value} settings, the critical feature here
is the name, \texttt{up} (which is at your discretion and could equally
be \texttt{upright} or \texttt{roman} or even \texttt{Fred} if you
so fancied). This ensures that the changed settings apply only to
this particular variant and do not `infect' the overall defaults.
To use this variant, all that is needed is to add the name, between
dots, to the \texttt{\textbackslash diff} command:
\begin{centred}
\verb`\[ \diff.up.yx \]` $\Longrightarrow$ \[ \diff.up.yx. \]
\end{centred}
\noindent Each variant derivative inherits all the default values
that it does not explicitly countermand. Thus a point of evaluation
is indicated by a vertical rule which is the \texttt{diffcoeff} default\footnote{\noindent But not the ISO recommendation.}:
\begin{centred}
\noindent \verb`\[ \diff.up.*{\frac{F(x)}{G(x)}}x[x=1] \]` $\Longrightarrow$
\[ \diff.up.*{\frac{F(x)}{G(x)}}x[x=1] \]
\end{centred}
\noindent Dot-delimited arguments must always be the \emph{first}
argument of the \texttt{\textbackslash diff} command, even preceding
an asterisk (star) as in this example.

As another example, suppose for the subscripted indication of variables
held constant in a partial derivative that you want to see what things
look like if the subscript is \emph{not }nudged in towards the right
parenthesis. In that case define a variant form 
\begin{lyxcode}
\textbackslash diffdef~\{~padrift~\}~\{~subscr-nudge~=~0~mu~\}
\end{lyxcode}
I have attached a name, \texttt{padrift},\texttt{ }to this setting,
\begin{centred}
\verb`\[ \diff.padrift.Fx[y] \]` $\Longrightarrow$ \[ \diff.padrift.Fx[y] \]
\end{centred}
\noindent since, to my eye, the subscript seems detached from the
expression it qualifies -- is it perhaps a typo? -- and `adrift
in a sea of whitespace'. This is to be compared with the default \verb`\[ \diffp Fx[y] \]`
$\Longrightarrow$ \[ \diffp Fx[y] \]

\noindent where the subscript is tucked in close to the right parenthesis
and is clearly connected to it and the expression it delimits.

Some might want to distinguish notationally a point of evaluation
for a partial derivative from variables held constant, perhaps using
a vertical rule for the former and (the default) parentheses for the
latter. It would suffice then to add to the preamble 
\begin{centred}
\verb`\diffdef { pvrule } { op-symbol = \partial }`
\end{centred}
This gives
\begin{centred}
\verb`\[ \diff.pvrule.{F(x,y)}x[x=1] \]`$\Longrightarrow$ \[ \diff.pvrule.{F(x,y)}x[x=1] \]
\end{centred}

\subsubsection{Text-style and script-style derivatives}

\label{subsec:Text-and-script-style}As noted earlier, the \texttt{diffcoeff}
package assumes that derivatives of fraction-like form will be used
in display-style expressions and that the slash form will be used
for inline use (text style). This is the usual practice. But if one
does want to use the fraction form in an inline expression, say \texttt{\textbackslash diffp
ST{[}V{]}} displaying as $\diffp ST[V]$, then some tweaking of settings
is necessary: the subscript is obviously too close to the right parenthesis
and, to my eye, there is too much `elbowroom' between the derivative
and the enclosing parentheses. Hence define

\noindent %
\noindent\begin{minipage}[t]{1\columnwidth}%
\begin{lyxcode}
\textbackslash diffdef~\{~ptxt~\}~

~~\{

~~~~op-symbol~~~~~~=~\textbackslash partial,

~~~~denom-term-sep~=~1~mu~~~~,

~~~~left-delim~~~~~=~\textbackslash left~(~,

~~~~right-delim~~~~=~\textbackslash right~),

~~~~elbowroom~~~~~~=~-2~mu~~~,

~~~~subscr-nudge~~~=~-3~mu~~~

~~\}
\end{lyxcode}
%
\end{minipage}

We can now write, for the same example, \texttt{\textbackslash diff.ptxt.ST{[}V{]}}
which displays as $\diff.ptxt.ST[V]$, where the subscript is better
positioned and there is a better fit between parentheses and derivative.

\subsubsection{Derivatives from D, \textbackslash delta, \textbackslash Delta}

\label{subsec:D-delta-Delta}In addition to \texttt{\textbackslash diff.p.},\texttt{
diffcoeff} has three other \emph{built-in} variant forms that are
commonly used: \texttt{\textbackslash diff.D.}, \texttt{\textbackslash diff.delta.},
and \texttt{\textbackslash diff.Delta.}, corresponding to derivatives
formed from $D$, $\delta$ and $\Delta$ respectively.

In fluid dynamics the \emph{material }or \emph{substantive} derivative
uses an uppercase $D$ in place of $d$. For example, the continuity
equation is,
\begin{centred}
\verb`\[ \diff.D.{\rho}t=\diffp\rho t + \mathbf{u\cdot}\nabla\rho \]`$\Longrightarrow$
\[ \diff.D.{\rho}t=\diffp\rho t + \mathbf{u\cdot}\nabla\rho \]
\end{centred}
\noindent where \texttt{\textbackslash diff.D.} produces the D-derivative.
If you want upright `D's, then the \texttt{ISO} package option will
produce that effect. Alternatively, 
\begin{lyxcode}
\textbackslash diffdef~\{~Up~\}~

~~\{

~~~~op-symbol~~~~=~\textbackslash mathrm\{D\},

~~~~op-order-sep~=~0~mu

~~\}
\end{lyxcode}
\noindent provides a variant with upright `D's.

In introductory calculus texts the simple $\delta$-derivative is
used. This is achieved with the \texttt{\textbackslash diff.delta.}
command:  \verb`\[ \diff.delta.yx \]` $\Longrightarrow$ \[ \diff.delta.yx. \]

\noindent This form also features in analytical mechanics (in the
Rogues' Gallery, the final example at (\ref{eq:eg6})).

Similarly, \texttt{\textbackslash diff.Delta.} forms a derivative
from $\Delta$:
\begin{centred}
\verb`$ \diff.Delta.y/x $` $\Longrightarrow$ $ \diff.Delta.y/x $
\end{centred}
\noindent where the slash form of the derivative is shown in this
instance.

Higher order forms of these derivatives, points of evaluation, appending
the differentiand with a star argument, all follow exactly as for
the `pure' \texttt{\textbackslash diff} command.

\paragraph{The commands \textbackslash Diff, \textbackslash diffd, \textbackslash Diffd}

For compatibility with version 1 of \texttt{diffcoeff}, the commands
\texttt{\textbackslash Diff}, \texttt{\textbackslash diffd} and
\texttt{\textbackslash Diffd} are available and also produce the
$D$, $\delta$ and $\Delta$ derivatives. Just as \texttt{\textbackslash diffp}
is identified with \texttt{\textbackslash diff.p.} for partial derivatives,
these commands are identified with \texttt{\textbackslash diff.D.},
\texttt{\textbackslash diff.delta.}, and \texttt{\textbackslash diff.Delta}
through the commands
\begin{lyxcode}
\textbackslash NewDocumentCommand~\textbackslash Diff~\{\}~~\{\textbackslash diff.D.\}

\textbackslash NewDocumentCommand~\textbackslash diffd~\{\}~\{\textbackslash diff.delta.\}

\textbackslash NewDocumentCommand~\textbackslash Diffd~\{\}~\{\textbackslash diff.Delta.\}
\end{lyxcode}
Unless one is using such variant forms frequently, it seems simpler
to remember that they are available as dot-delimited arguments to
the \texttt{\textbackslash diff} command, using the obvious name
in each case, rather than having to remember the precise camel-case
form of name of the \texttt{\textbackslash Diff}, \texttt{\textbackslash diffd}
and \texttt{\textbackslash Diffd} commands.

\subsection{The \texttt{.def} file}

\label{subsec:The-.def-file}This mechanism of variant formation is
how I have been able to illustrate in the Rogues' Gallery, �\ref{sec:Rogues'-gallery},
the wide variety of different usages culled from the literature. But
the thought arises: if a variant is to be used only once or twice,
isn't this a lot of bother? Why not just construct the variant derivative
`by hand' out of \texttt{\textbackslash frac} and \texttt{\textbackslash mskip}
for example? The reason for making such definitions is that they can
be transferred from document to document. For instance, definitions
placed in the preamble can be copied to the preamble of another document.

But that is hardly optimal. Instead, \texttt{diffcoeff} allows such
definitions to be placed in a text file with the the extension \texttt{.def}
and a name of your choosing. For the present document the file is
called \texttt{diffcoeff-doc.def} and should have been placed alongside
the \verb`.sty` file in your \TeX{} distribution. The contents of
this file can be seen at �To use these definitions, the \texttt{diffcoeff}
package is called with the command 
\begin{lyxcode}
\textbackslash usepackage{[}def-file=diffcoeff-doc{]}\{diffcoeff\}{[}=v4{]}
\end{lyxcode}
But even this process still means copying a definition file from directory
to directory as one works on different documents. The solution is
to make a definition file available for \emph{all} documents and the
way to do that is by placing it in the texmf tree, preferably not
the one created by your \TeX{} distribution (perhaps MiKTeX or TexLive),
but your own \emph{personal} texmf tree.

\texttt{\vspace{2ex}
 }

\noindent %
\noindent\fbox{\begin{minipage}[t]{1\linewidth - 2\fboxsep - 2\fboxrule}%
\textbf{Personal texmf tree? }

This is a directory for `waifs and strays' of the \TeX{} system that
are not included in standard distributions like MiK\TeX{} or \TeX Live.
For instance, it is the place for personal packages designed for your
own particular circumstances or preferences, and is structured like
the standard MiK\TeX{} or \TeX Live hierarchy but placed in another
location so that there is no chance of its being overwritten when
MiK\TeX{} or \TeX Live are updated. However, those distributions need
to be alerted to its existence. For MiK\TeX , open the MiK\TeX{} console,
click on \textsf{Settings} and then the \textsf{Directories} tab.
Click the \textsf{+} button and navigate to your personal texmf tree
to add it to the MiK\TeX{} search path. Having added it, you will then
need to refresh the filename database by clicking on the \textsf{Tasks}
menu and selecting the obvious entry. I am not familiar with \TeX Live
but presume an analogous process will apply there.%
\end{minipage}}

\medskip{}

\noindent Provided your \LaTeX{} distribution knows about your personal
texmf tree, then a \texttt{.def} file placed within it, will be accessible
to all documents.

\subsubsection{Structure of the \texttt{.def} file}

The best way to see what a \texttt{.def} file looks like is to view
\texttt{diffcoeff-doc.def }in a text editor.

If you want your variant definitions to use defaults different from
those supplied with the \texttt{diffcoeff} package, then the first
definition in the \texttt{.def} file should be the one setting the
new defaults, with an \emph{empty} first argument to the \texttt{\textbackslash diffdef}
command: 
\begin{lyxcode}
\textbackslash diffdef~\{\}~~~~

~~\{~~~~~~

~~~~key-1~=~value-1,~~~~

~~~~key-2~=~value-2,

~~~~...

~~~~key-n~=~value-n

~~\}
\end{lyxcode}
The key-value list is a comma-separated list; hence the last entry
doesn't need to end with a comma. Nudge and separation values need
to include the unit, \texttt{mu}; a numerical value alone will result
in error.\texttt{ }Because a \texttt{.def} file\texttt{ }is a \LaTeX{}
file, comments need to start with a \texttt{\%} character.

\subsubsection{\texttt{diffcoeff.def}}

Note that if the \texttt{diffcoeff} package is invoked without an
explicit \texttt{def-file= <filename>} option statement, as here, 
\begin{lyxcode}
\textbackslash usepackage\{diffcoeff\}{[}=v4{]}
\end{lyxcode}
then it will search in the texmf tree (the \LaTeX{} distribution's
and your personal one) and the document directory for a file \texttt{diffcoeff.def}
and if found will load that. This file should contain definitions
of those variants you are likely to use in multiple documents. In
my personal texmf tree (which I've put at \texttt{E:\textbackslash texmf\textbackslash}
on a Windows machine) the file \texttt{diffcoeff.def} is located in
the directory \texttt{E:\textbackslash texmf\textbackslash tex\textbackslash latex\textbackslash diffcoeff\textbackslash}.
(The backslashes are replaced by forward slashes on linux machines.)

Variants likely to be of value only to a specific document should
be added to the preamble of that document. Alternatively, they could
be added to \texttt{diffcoeff.def} but that added-to file saved to
the document directory under a \emph{different} name -- e.g. I've
saved the variants required for the present document under the name
\texttt{diffcoeff-doc.def}. Many of these variants were created solely
to illustrate points in the present document and I have no intention
of using them in my own work. Consequently my \texttt{diffcoeff.def}
file is smaller, containing only a selection from \texttt{diffcoeff-doc.def}.

\section{Differentials in integrals, etc.}

Forms like $dx$ occur not only in derivatives but also in other contexts,\footnote{I thank \noun{Sergio Callegari} for drawing my attention to this rather
obvious point, and its omission from previous versions of \texttt{diffcoeff}.} for example, the expression for a total differential like 
\[
\dl P=\diffp Px\dl x+\diffp Py\dl y+\diffp Pz\dl z,
\]
or an integral like $\int\sin x\dl2x$, or a multi-variable integral
like 
\[
\iiintop_{-\infty}^{\infty}V(x,y,z)\dl3x\dl2y\dl2z.
\]
Surely we want the `d's in these expressions to correspond to their
form (upright or math italic) in derivatives? If, for instance, the
\texttt{ISO} package option has been set, one doesn't want to be writing
\texttt{\textbackslash mathrm\{d\}} in every (or indeed any) integral.
To this end, \texttt{diffcoeff} provides a command \texttt{\textbackslash dl}
to write the `d' in a differential in a manner consistent with the
default form used in derivatives. In the present document, the default
form is math-italic and so
\begin{centred}
\verb`$ \dl x $` $\Longrightarrow$ $ \dl x. $
\end{centred}
\noindent It is also possible\footnote{\noindent From version 3.1 of \texttt{diffcoeff}; version 3.0 of the
package produced a \LaTeX{} error.} to use the command before, for instance, \texttt{\textbackslash mathbf\{x\}}
or \texttt{\textbackslash vec\{x\}}:
\begin{centred}
\verb`$ \dl \vec{x} $` $\Longrightarrow$ $ \dl \vec{x} $,~~~~\verb`$ \dl \mathbf{x} $`
$\Longrightarrow$$ \dl \mathbf{x}. $
\end{centred}

\subsection{Options}

There are two options available with the differential command.

The first is the dotted name option discussed in �\ref{sec:Changing-defaults}.
Thus, to illustrate a differential with an upright `d', write
\begin{centred}
\verb`$ \dl.up.x $` $\Longrightarrow$ $ \dl.up.x. $
\end{centred}
\noindent This presumes that a variant derivative with upright `d's
has been defined and the definition is available, either in the preamble
or in an accessible \texttt{.def} file, as discussed in the previous
section. Alternatively, if you have defined your default derivative
to use upright `d's (perhaps with the \texttt{ISO} option) then
\texttt{\textbackslash dl} alone will suffice to produce an upright
`d'.

\subsubsection{Partial differential}

\label{subsec:Partial-differential}Since the variant \texttt{\textbackslash diff.p.}
is defined in \texttt{diffcoeff4.sty} itself, \texttt{\textbackslash dl.p.}
is always available and at 6 keystrokes offers a slightly shorter
way of writing \texttt{\textbackslash partial} (8 keystrokes). However,
there are sufficient contexts where expressions like \verb`\partial_x`
are used, perhaps as a shortcut for a partial derivative, for it to
be worthwhile to define an appropriate command for this variant. From
version 4.0, \texttt{diffcoeff} therefore provides \verb`\dlp`, defined
by
\begin{lyxcode}
\textbackslash NewDocumentCommand~\textbackslash dlp~\{\}~\{~\textbackslash dl.p.~\}
\end{lyxcode}
Thus, for instance, \verb`$ \dlp_x $` $\Longrightarrow$ $ \dlp_x $.

\subsubsection{Spacing}

This second option inserts spacing before the `d'. If \texttt{\textbackslash dl}
is followed by a digit (0, 1, 2, ..., 9) it will insert a horizontal
space of that number of mu before the `d'; (\texttt{\textbackslash dl0x}
has the same effect as \texttt{\textbackslash dl x}.)\texttt{ }Thus,
an alternative way of writing an example in Chapter~18 of \emph{The
\TeX book} is
\begin{centred}
\verb`$ \dl x\dl3y=r\dl3r\dl3\theta $` $\Longrightarrow$ $ \dl x\dl3y=r\dl3r\dl3\theta. $
\end{centred}
\noindent To my eye this is too much space; I prefer
\begin{centred}
\noindent \verb`$ \dl x\dl2y=r\dl2r\dl2\theta $` $\Longrightarrow$
$ \dl x\dl2y=r\dl2r\dl2\theta. $
\end{centred}
\noindent I used \texttt{\textbackslash dl3x\textbackslash dl2y\textbackslash dl2z}
when writing the triple integral above, \emph{no} extra spacing when
writing the total differential expression, since the differentials
are already distinct from the preceding fraction forms $\diff.ptxt.Px$
etc., and \texttt{\textbackslash dl2x} when writing $\int\sin x\dl2x$.

Note that only \emph{one} digit can be used. If two are used, as here,
\verb`$ \dl20x $` $\Longrightarrow$ $ \dl20x $, the effect is unlikely
to be what is wanted.

From version 3.2 of \texttt{diffcoeff} it is possible to also add
\emph{negative} space before the differential, which might be useful
in special contexts, perhaps to construct a symbol:
\begin{centred}
\verb`$ /\dl-9x $` $\Longrightarrow$ $ /\dl-9x. $
\end{centred}
\noindent For negative space before the differential, add a minus
sign before the (single) digit.

The spacing digit option \emph{follows} the dot-delimited name option.
For example, earlier, in the denominator of a mixed partial derivative,
I have used what is effectively \texttt{\textbackslash dlp x\textbackslash dlp2y\textbackslash dlp2z},
and the variation to that spacing when a higher-order differentiation
occurs:
\begin{centred}
\verb`$ \dlp x^2\dlp1y\dlp2z $` $\Longrightarrow$ $ \dlp x^2\dlp1y\dlp2z. $
\end{centred}

\subsection{Rationale}

Version 1 of the \texttt{diffcoeff} package arose from a need to simplify
the parsing of differential coefficients for another program I was
working on which was struggling to `read' all the possible permutations
of \texttt{\textbackslash frac} or \texttt{\textbackslash tfrac}
or \texttt{\textbackslash dfrac} or slash forms of the derivative,
of \texttt{d} or \texttt{\textbackslash mathrm\{d\}} or \texttt{\textbackslash partial}
or \texttt{D} or \texttt{\textbackslash mathrm\{D\}} or \texttt{\textbackslash nabla},\texttt{
}and of points of evaluation delimited by vertical rules or parentheses.\texttt{ }Although
regular expressions coped with most of these cases, it was \emph{messy}.

There are other packages which have commands for the derivative (e.g.,
\texttt{bropd}, \texttt{commath},\texttt{ esdiff}, \texttt{physymb})
but none quite gave what I wanted -- although they probably cope
with most users' needs. \texttt{esdiff} came closest to what I was
seeking but failed when it came to combining algebraic and numeric
orders of differentation in a mixed partial derivative (and made heavier
use of braces than I liked in that case too).

\subsubsection{\texttt{diffcoeff.sty}}

I have tried to make using \texttt{diffcoeff} intuitive. Looking at
the other packages mentioned, writing something like \texttt{\textbackslash diff{[}n{]}\{f\}\{x\}}
(which can be trimmed to \texttt{\textbackslash diff{[}n{]}fx} for
single-token arguments) seems `natural' -- only \texttt{physymb}
deviates from the pattern. 
\begin{itemize}
\item It seems consistent with this pattern to use a comma list for specifying
the orders of differentiation of the variables in a higher order mixed
partial derivative (and its suppression when all are of order 1) 
\item Having specified the orders, surely the program itself should calculate
the overall order? \texttt{esdiff} does this for numerical orders;
\texttt{diffcoeff} does this for both numeric and algebraic orders, 
\end{itemize}
\begin{centred}
\verb`\[ \diffp[m-(k+1),m+(k-1)]{F(x,y,z)}{x,y,z} \]`
\end{centred}
\begin{example}
$\Longrightarrow$ \[ \diffp[m-(k+1),m+(k-1)]{F(x,y,z)}{x,y,z} \]
\end{example}

\begin{itemize}
\item and where it fails either to calculate at all or to present the result
in a preferred form, offers the order-override option: 
\end{itemize}
\begin{centred}
\verb`\[ \diffp[m+(k+1),m+(k-1)][2(m+k+1)]{F(x,y,z,w)}{x,y,z,w} \]`
\end{centred}
\begin{example}
$\Longrightarrow$ \[ \diffp[m+(k+1),m+(k-1)][2(m+k+1)]{F(x,y,z,w)}{x,y,z,w} \]
\end{example}

\begin{itemize}
\item I wished to avoid the unnecessary writing of superscripts, subscripts
and brace pairs. In the examples just given, no superscript tokens
\texttt{\textasciicircum} are written by the user despite the higher-order
differentiation in $x$ and $y$, and only the two inescapable brace
pairs are required. 
\item The use of a comma list for the second mandatory argument in a partial
derivative -- the list of variables -- makes differentiations in
super- or subscripted symbols (as occurs prolifically in tensor calculus)
easier to both write and read by avoiding unnecessary brace pairs. 
\end{itemize}
\begin{centred}
\verb`\[ \diffp{A_i}{ x^j,x^k } \]` $\Longrightarrow$ \[ \diffp{A_i}{ x^j,x^k } \]
\end{centred}
\begin{itemize}
\item Should a point of evaluation or variables held constant be considered
part of the derivative? Thermodynamic usage was decisive here. The
partial derivative alone is ambiguous -- the parentheses and subscript
are essential to understand what is being stated: 
\[
\diffp ST[V]
\]
Hence provision for these extra elements was included in the derivative
commands. 
\item Given the position of the subscripted symbol in the displayed derivative,
it's positioning as the \emph{final} argument in the derivative command
feels inevitable. 
\item Version 1 of \texttt{diffcoeff} used braces for this argument to avoid
any possible confusion with a following mathematical expression. That
use of braces is now deprecated in \texttt{xparse}, has been deprecated
in \texttt{diffcoeff} since version 2, and is no longer compatible
with version 4. Later versions of \texttt{diffcoeff} use square brackets,
conforming with familiar \LaTeX{} practice. The only special remembering
needed is avoidance of a space before the argument -- and if it does
slip in, it won't cause a \LaTeX{} error. It will be treated as part
of a following mathematical expression and displayed as such. 
\item The star option also prompted the question: is it needed? After all,
one can always leave the first mandatory argument empty and append
the differentiand `by hand'. But once the provision for points of
evaluation or variables held constant was incorporated into the derivative
commands, the star option became the simplest way of handling appended
differentiands since the parentheses for a variable held constant
must wrap around the differential operator \emph{and} differentiand.
Once available, it provides a simple way of switching between (and
comparing) the appearance of differentiand-in-the-numerator and differentiand-appended. 
\item The slash option was added to the derivative commands after seeing
how widely such forms are used in texts at all levels. The placement
of the slash, between the two mandatory arguments, seems more-or-less
self-evident. 
\item A later option added to \texttt{\textbackslash diff} (and not present
in version 1) was the dot-delimited name option. Once \texttt{xtemplate}
was used as the basis of the package this seemed the most straightforward
way of making available, ready to hand, the wealth of variants that
\texttt{xtemplate} makes possible. (It's just a pity that the second
dot is needed, and a single-dot naming scheme can't be used, but \texttt{xparse}
forces my hand here.) 
\item Having added the dot-delimited name option, the use of a \texttt{def}
file to store variants or preferred defaults is more-or-less forced,
otherwise one is faced with making these definitions anew for each
new document (or locating a previous document and copying from that
to the new one). 
\item To handle possible differences between display-style and text-style
(and script-style) derivatives (see �\ref{subsec:Text-and-script-style})
I considered using \TeX 's \texttt{\textbackslash mathchoice} command.
This command takes four arguments, corresponding to display-, text-,
script- and scriptscript-styles and would require forming four derivatives
each time a derivative is used, `just in case'. In fact fraction-form
derivatives are used overwhelmingly in display-style expressions,
the slash form being used for inline use. Given the ease of defining
a fraction-form variant for text-style use, and the rareness of such
use, employing variants seemed the way to go. It is the one adopted
and avoids the computational burden associated with the use of \texttt{\textbackslash mathchoice}. 
\item After version 2 of the package appeared on CTAN, it was pointed out
to me that there was an issue of consistency between the form of `d'
used in a derivative (upright or math-italic) and the form used in
an integral. I had overlooked this matter completely and in version
3 of the package remedied the omission with the differential command
\texttt{\textbackslash dl}. A spacing option for \texttt{\textbackslash dl}
was almost inevitable. From version 4 the partial analogue \texttt{\textbackslash dlp}
has been added.
\item Space before the differentiand was requested by a user. Once considered
it became clear that there are (at least) two ways of thinking of
a derivative: as $\diff y/x$, a ratio of differentials where $\dl y$
is a unit and it makes no sense to insert space between the `d'
and the `y', and as
\[
\diff!{F(x)}x
\]
where the function $F(x)$ is being operated on by $\diff{}/x$ and
it is natural to insert space between the `d' and the `F'. With
that realisation came the need for a simple switch-like package option
(\verb`spaced`) to turn spacing on or off, and a switch-like argument
(\verb`!`) to countermand the package option in exceptional cases
(version 4).
\end{itemize}

\section{Reference}

\subsection{Commands}
\begin{description}
\item [{\texttt{\textbackslash diff}}] arguments (all optional unless
otherwise indicated): 
\begin{enumerate}
\item .\emph{name}. for the given settings 
\item {*} append-differentiand switch 
\item {[}\emph{order}{]} or {[}\emph{comma-list of orders}{]} of differentiation 
\item {[}\emph{order-override}{]} 
\item ! countermand before-differentiand spacing of the \verb`spaced` package
option
\item \{\emph{differentiand}\} (mandatory) 
\item / slash-form switch 
\item \{\emph{comma list of differentiation variables}\} (mandatory) 
\item {[}\emph{point of evaluation/variables held constant}{]} 
\end{enumerate}
\item [{\texttt{\textbackslash diffdef}}] arguments (all mandatory): 
\begin{enumerate}
\item \{\emph{name}\} 
\item \{\emph{key=value comma list}\} 
\end{enumerate}
\item [{\texttt{\textbackslash diffp}}] \texttt{= \textbackslash diff.p.}
\item [{\texttt{\textbackslash dl}}] arguments (all optional): 
\begin{enumerate}
\item .\emph{name}. (as for \texttt{\textbackslash diff}) 
\item \emph{minus sign} (optional, use only if negative spacing before the
`d' is wanted)
\item \emph{digit} (insert spacing of this number of mu before the `d') 
\end{enumerate}
\item [{\texttt{\textbackslash dpl}}] \texttt{= \textbackslash dl.p.}
\item [{\texttt{\textbackslash negmu}}] insert a $-1$ mu space 
\item [{\texttt{\textbackslash nilmu}}] insert a $0$ mu space 
\item [{\texttt{\textbackslash onemu}}] insert a $1$ mu space 
\item [{\texttt{\textbackslash twomu}}] insert a $2$ mu space
\end{description}
Deprecated (why try remembering the special names when the forms on
the right are self-explanatory): 
\begin{description}
\item [{\texttt{\textbackslash Diff}}] preferred form: \texttt{\textbackslash diff.D.} 
\item [{\texttt{\textbackslash diffd}}] preferred form \texttt{\textbackslash diff.delta.}
\item [{\texttt{\textbackslash Diffd}}] preferred form\texttt{ \textbackslash diff.Delta.}
\end{description}

\subsection{The file \texttt{diffcoeff-doc.def}}

\label{sec:diffcoeff-doc.def}The file \verb`diffcoeff-doc.def` contains
the following definitions:
\begin{verbatim}
%%%%%%%%%% ordinary %%%%%%%%%%
% ( ) for pt of eval.
\diffdef { paren } 
  {
    left-delim   = \left ( ,
    right-delim  = \right ),
    subscr-nudge = -6 mu
  }

% upright
\diffdef { up } 
  {
    op-symbol      = \mathrm{d},
    op-order-sep   = 0 mu 
  }

\diffdef { Up } 
  {
    op-symbol      = \mathrm{D},
    op-order-sep   = 0 mu
  }

% wrap long vars: (d longvar)
\diffdef { wrapall } {long-var-wrap = (dv)}

% nabla in numerator, d in denominator
\diffdef { nabla } 
  {
    op-symbol     = \nabla,diffcoeff-doc.def
    op-symbol-alt = d
  }

% align op left; no italic nudge 
\diffdef { left } 
  {
    *-op-left      = true,
    *-italic-nudge = 0 mu
  }

% slash-form: vrule for pt of eval
\diffdef { svrule } 
  {
    /-left-delim  =  ,
    /-right-delim = |
  }

%%%%%%%%%% partial %%%%%%%%%%
% partial, vrule for pt of eval.
\diffdef { pvrule } { op-symbol = \partial }

% partial, no subscr nudge (adrift in a sea of white space)
\diffdef { padrift } 
  { 
    op-symbol    = \partial,
    left-delim   = \left ( ,
    right-delim  = \right ),
    subscr-nudge = 0 mu
  }

% partial, [ ] for pt of eval.
\diffdef { psqbra } 
  {
    op-symbol    = \partial,
    left-delim   = \left [ ,
    right-delim  = \right ],
    subscr-nudge = 0 mu
  }

% Delta as partial deriv.
\diffdef { pDelta } 
  {
    op-symbol      = \Delta  ,
    op-order-sep   = 0 mu    ,
    left-delim     = \left ( ,
    right-delim    = \right ),
    subscr-nudge   = -6 mu
  }

% align op left; italic nudge
\diffdef { pleft } 
  {
    op-symbol      = \partial,
    *-op-left      = true    ,
    *-italic-nudge = 3 mu    ,
    left-delim     = \left ( ,
    right-delim    = \right ),
    subscr-nudge   = -3 mu
  }

% partial, textstyle
\diffdef { ptxt } 
  {
    op-symbol      = \partial,
    denom-term-sep = 1 mu    ,
    left-delim     = \left ( ,
    right-delim    = \right ),
    elbowroom      = -2 mu   ,
    subscr-nudge   = -3 mu 
  }

% partial, wide sep of terms
\diffdef { pwide } 
  {
    op-symbol        = \partial,
    denom-term-sep   = 3 mu    ,
    /-denom-term-sep = 3 mu    ,
    left-delim       = \left ( ,
    right-delim      = \right ),
    subscr-nudge     = -6 mu 
  }
 
 %%%%%%%%%% differential %%%%%%%%%%
 \diffdef { up } 
  { op-symbol = \mathrm{d} }
\end{verbatim}


\subsection{Version comparison}

\label{sec:Version-comparison}Unlike version 1, version 2 and later
are built on the the \texttt{xtemplate} package which makes certain
facilities available which it would be silly not to exploit. Hence
the coding in the later versions is completely different and there
are consequences.
\begin{enumerate}
\item From version 2.0
\begin{enumerate}
\item The \texttt{\textbackslash diffset} command, formerly used to tweak
the display of derivatives, has been superseded by the \texttt{\textbackslash diffdef}
command. \texttt{\textbackslash diffset} now sends a message warning
of its obsolescence to the terminal and \LaTeX{} log but is otherwise
functionless. It should not interfere with the compilation of a document
but any intended fine-tuning of the display by means of the \texttt{\textbackslash diffset}
command\texttt{ }will not eventuate. The warning message is: \texttt{Obsolete
command: \textbackslash diffset has been superseded by the \textbackslash diffdef
command.} \texttt{See the diffcoeff} \texttt{doc\-umentation for
further information.} The \texttt{\textbackslash diffdef} command
is discussed in �\ref{subsec:diffdef}; 
\item The optional trailing argument used to indicate a point of evaluation
or variables held constant is now delimited by square brackets, \texttt{{[}}
and \texttt{{]}}, as other optional arguments are. For compatibility
with version 1, versions 2 and 3 still accepted braces to delimit
this argument but from version 4 of \texttt{diffcoeff} only the square-bracket
delimited argument is accepted. (The use of braces around \emph{optional}
arguments while once accepted is now deprecated in \texttt{xparse}
on which \texttt{diffcoeff} depends);
\item The commands \texttt{\textbackslash Diff}, \texttt{\textbackslash diffd}
and \texttt{\textbackslash Diffd} used to construct derivatives from
$D$, $\delta$ and $\Delta$ in version 1, are still available but
deprecated. A new optional argument in the \texttt{\textbackslash diff}
command offers these and a host of other possibilities and is now
the preferred method of forming such variants; see �\ref{subsec:D-delta-Delta}.
\end{enumerate}
\item Version 3.0
\begin{enumerate}
\item adds a command, \texttt{\textbackslash dl} (from \emph{d}ifferentia\emph{l})
to write differentials like $dx$ that occur in integrals and in other
contexts in a manner consistent with the form used in derivatives.
After all, if one is using upright `d's in derivatives, similarly
upright `d's should occur in these other contexts;\footnote{This rather obvious lack in version 2 was pointed out to me by Sergio
Callegari.}
\item provides some simple spacing commands that can be useful for tweaking
standard spacing.
\end{enumerate}
\item Version 3.1 enables the differential command to be used before forms
like \texttt{\textbackslash vec\{x\}} (an overlooked possibility
causing an error in earlier versions).
\item Version 3.2
\begin{enumerate}
\item allows negative spacing before the differential command \texttt{\textbackslash dl};
\item fixes a bug in which an ordinary derivative as the differentiand of
a partial derivative displayed as a partial derivativ\texttt{diffcoeff-doc.def}e.
It now displays, as it should, as an ordinary derivative. 
\end{enumerate}
\item Version 4.0 
\begin{enumerate}
\item enables the insertion of a small space before the differentiand, either
as the default behaviour (package option \verb`spaced`) or at explicit
request (argument \verb`!` of the \verb`\diff` command); see �\ref{subsec:Spacing-before-derivand}. 
\item prevents the ligature $df$ that marred previous versions; this is
now rendered $\dl f$.
\item offers the document command \verb`\dlp` for the \emph{partial} differential;
see �\ref{subsec:Partial-differential};
\item no longer accepts \emph{the braced form} of the\emph{ }trailing optional
argument specifying a point of evaluation or (for partial derivatives)
variables held constant. This was a relic from version 1 of \texttt{diffcoeff},
and has been deprecated since version 2.
\end{enumerate}
\item Version 4.1
\begin{enumerate}
\item changes the code for the \verb`long-var-wrap` setting to avoid a
problem arising from a change in \verb`xtemplate`; 
\item improves spacing between operator and variable in \emph{denominators}
of differential coefficients (e.g. prevents ligatures).
\end{enumerate}
\item Version 4.2
\begin{enumerate}
\item see the announcement at the start of this document.
\end{enumerate}
\end{enumerate}

\end{document}
